\chapter{Sistem nonlinear} \label{nlin:chapter}

%%%%%%%%%%%%%%%%%%%%%%%%%%%%%%%%%%%%%%%%%%%%%%%%%%%%%%%%%%%%%%%%%%%%%%%%%%%%%%

\section{Linearisasi, titik kritis, dan kesetimbangan}
\label{linearization:section}

%mbxINTROSUBSECTION

\sectionnotes{1 pertemuan\EPref{, \S6.1--\S6.2 dalam \cite{EP}}\BDref{,
\S9.2--\S9.3 dalam \cite{BD}}}

%\subsection{Persamaan nonlinear}

Selain beberapa selingan singkat dalam \chapterref{fo:chapter},
kita terutama membahas persamaan
linear. Persamaan linear memadai dalam banyak penerapan, tetapi dalam kenyataan,
kebanyakan fenomena memerlukan persamaan nonlinear. Namun, persamaan nonlinear
terkenal jauh lebih sulit dipahami daripada persamaan linear, dan
banyak fenomena baru yang aneh muncul ketika kita mengizinkan persamaan kita
bersifat nonlinear.
Jangan khawatir, waktu yang kita gunakan untuk mempelajari persamaan linear tidak sia-sia.
Persamaan nonlinear sering kali dapat dihampiri dengan persamaan linear jika kita hanya memerlukan
solusi secara \myquote{lokal}---hanya untuk jangka waktu singkat atau
hanya untuk parameter tertentu. Pemahaman terhadap persamaan linear juga dapat
memberi kita pemahaman kualitatif tentang masalah nonlinear yang lebih umum.
Gagasannya serupa dengan yang Anda lakukan dalam kalkulus ketika mencoba
menghampiri suatu fungsi dengan garis yang memiliki kemiringan yang tepat.

\begin{mywrapfigsimp}{1.45in}{1.75in}
\noindent
\inputpdft{mv-pend}{%
%note that this diagram appears more than once in the book
Diagram sebuah bandul. Sebuah bandul diperlihatkan terpancang pada satu titik di atas
dan bergerak sepanjang busur lingkaran. Ujung bawah bandul diberi label m, dan
sudut antara bandul dan garis vertikal diberi label theta.
Panjang bandul diberi label L.}
\end{mywrapfigsimp}
Dalam \sectionref{sec:mv} kita meninjau bandul dengan
panjang $L$. Tujuannya adalah mencari sudut $\theta(t)$ sebagai
fungsi waktu $t$. Persamaan untuk susunan ini adalah
persamaan nonlinear
\begin{equation*}
\theta'' + \frac{g}{L} \sin \theta = 0 .
\end{equation*}
Alih-alih menyelesaikan persamaan ini, kita menyelesaikan persamaan linear yang
jauh lebih mudah
\begin{equation*}
\theta'' + \frac{g}{L} \theta = 0 .
\end{equation*}
Meskipun solusi persamaan linear bukan persis yang kita
cari, solusi itu cukup dekat dengan solusi semula, selama
sudut $\theta$ kecil dan selang waktunya singkat.

Anda mungkin bertanya: Mengapa kita tidak langsung menyelesaikan masalah nonlinear itu? Nah,
persamaan yang dimaksud mungkin sangat sulit, tidak praktis, atau mustahil diselesaikan
secara analitik. Kita bahkan mungkin
tidak tertarik pada solusi sebenarnya.
Kita mungkin hanya tertarik
pada gambaran kualitatif mengenai perilaku solusi. Misalnya,
%we may be interested in
apa yang terjadi ketika waktu menuju tak hingga?
%In the case
%of the pendulum we found that it oscillates and we can even approximate
%the period well if the swings are small.
%In other words, why do more work, when we can do less.
%The exact solution, even if found, might be harder to analyze.


\subsection{Sistem otonom dan analisis bidang fase}

Kita membatasi perhatian pada sistem otonom berdimensi dua
\begin{equation*}
x' = f(x,y) , \qquad y' = g(x,y) ,
\end{equation*}
dengan $f(x,y)$ dan $g(x,y)$ merupakan fungsi dua variabel, dan
turunannya diambil terhadap waktu $t$. Solusinya adalah
fungsi $x(t)$ dan $y(t)$ sedemikian sehingga
\begin{equation*}
x'(t) = f\bigl(x(t),y(t)\bigr), \qquad
y'(t) = g\bigl(x(t),y(t)\bigr) .
\end{equation*}
Cara kita menganalisis sistem ini sangat mirip dengan
\sectionref{auteq:section}, tempat kita mempelajari satu persamaan otonom. Gagasan
dalam dua dimensi sama, tetapi perilakunya dapat
jauh lebih rumit.
%We will do the same sort of analysis.  We will
%look for the \emph{critical points} of the system and then we will analyze
%what happens when time goes to infinity.

Mungkin yang terbaik adalah memandang sistem persamaan tersebut sebagai satu persamaan vektor
\begin{equation} \label{eq:nlinautn2}
\begin{bmatrix} x \\ y \end{bmatrix} ' =
\begin{bmatrix} f(x,y) \\ g(x,y) \end{bmatrix} .
\end{equation}
Seperti dalam \sectionref{sec:introtosys}, kita menggambar
\emph{\myindex{potret fase}} (atau \emph{\myindex{diagram fase}}),
tempat setiap titik $(x,y)$ bersesuaian dengan keadaan tertentu dari sistem.
Kita menggambar \emph{\myindex{medan vektor}}
yang pada setiap
titik $(x,y)$ diberikan oleh vektor
$\left[ \begin{smallmatrix} f(x,y) \\ g(x,y) \end{smallmatrix} \right]$.
Dan seperti sebelumnya, jika kita menemukan solusi, kita menggambar lintasannya
dengan memplot semua titik $\bigl(x(t),y(t)\bigr)$ untuk suatu rentang $t$.

\begin{example} \label{example:nlin-1b-example}
Tinjau persamaan orde kedua $x''=-x+x^2$.
Tuliskan persamaan ini sebagai sistem nonlinear orde pertama
\begin{equation*}
x' = y , \qquad y' = -x+x^2 .
\end{equation*}
Potret fase beserta beberapa lintasannya digambar dalam
\figurevref{fig:nlin-1b}.
\begin{myfig}
\capstart
\diffyincludegraphics{width=3in}{width=4.5in}{nlin-1b}{%
Plot medan arah (kisi panah) pada bidang xy.
Di tengah, panah-panah mengitari titik kritis di
titik asal searah jarum jam. Kemudian,
pada titik (1,0), medan arahnya berupa titik pelana, yaitu,
panah-panah sepanjang kurva miring turun secara diagonal menuju titik kritis,
sedangkan pada kurva miring naik secara diagonal menjauhi titik kritis.
Di sebelah kiri titik kritis, panah-panah kira-kira mengarah ke bawah; di atas titik kritis,
panah-panah kira-kira mengarah ke kanan;
di sebelah kanan titik kritis, panah-panah kira-kira mengarah ke atas;
dan di bawah titik kritis, panah-panah kira-kira mengarah ke kiri.
Panah-panah tampak semakin panjang ketika semakin jauh dari titik kritis.
Beberapa lintasan juga digambar. Di sekitar titik kritis
(0,0) terdapat beberapa lintasan berbentuk telur dengan ujung runcing ke arah
x positif. Lalu ada satu lintasan yang tampak
melewati titik kritis (1,0) sambil membentuk gelung, walaupun untuk benar-benar
membentuk gelung, lintasan harus melawan arah panah setelah melewati
titik kritis. Kemudian terdapat beberapa lintasan lain;
beberapa lintasan berawal dari kanan bawah, naik secara diagonal,
tetapi kemudian berbelok ke kiri hingga mencapai separuh kiri gambar
dengan tetap berada di bawah sumbu x.
Kemudian lintasan-lintasan itu berbelok ke atas hingga
berada di separuh atas gambar, lalu berbelok ke kanan hingga mencapai
separuh kanan gambar, tempat lintasan itu berbelok naik secara diagonal dan meninggalkan
gambar. Sepasang lintasan lain berawal di bawah sumbu x pada
sisi kanan gambar, dekat dengan sumbu x, lalu sekadar berbelok
ke kanan hingga meninggalkan gambar, tanpa pernah melewati x sama dengan 1.}
\caption{Potret fase beserta beberapa lintasan dari
$x' = y$, $y' = -x+x^2$. \label{fig:nlin-1b}}
\end{myfig}

Dari potret fase tampak jelas bahwa bahkan sistem sederhana ini memiliki
perilaku yang cukup rumit. Beberapa lintasan terus berosilasi mengitari
titik asal, dan beberapa lainnya bergerak menuju tak hingga. Kita akan sering
kembali ke contoh ini dan menganalisisnya secara lengkap dalam bagian ini (dan bagian berikutnya).
\end{example}

Jika kita memperbesar diagram di dekat suatu titik tempat
$\left[ \begin{smallmatrix} f(x,y) \\ g(x,y) \end{smallmatrix} \right]$
tidak nol, maka panah-panah di sekitarnya pada umumnya menunjuk ke arah yang pada dasarnya sama
dan memiliki besar yang pada dasarnya sama.
Dengan kata lain, perilaku di dekat titik seperti itu tidak terlalu menarik.
Tentu saja kita mengasumsikan bahwa $f(x,y)$ dan $g(x,y)$ kontinu.

Mari kita pusatkan perhatian pada titik-titik dalam diagram fase
di atas, tempat lintasan-lintasan
tampak berawal, berakhir, atau mengitari sesuatu. Kita melihat dua titik demikian:
$(0,0)$ dan $(1,0)$. Lintasan-lintasan tampak mengitari titik $(0,0)$,
dan tampak masuk ke atau keluar dari titik $(1,0)$.
%
Titik-titik ini tepat merupakan titik tempat turunan $x$
dan $y$ sama-sama nol. Kita definisikan \emph{titik kritis}\index{titik kritis}
sebagai titik $(x,y)$ sedemikian sehingga
\begin{equation*} 
\begin{bmatrix} f(x,y) \\ g(x,y) \end{bmatrix} = \vec{0} .
\avoidbreak
\end{equation*}
Dengan kata lain, inilah titik-titik tempat $f(x,y)=0$ dan $g(x,y)=0$.

Titik kritis adalah tempat perilaku sistem, dalam arti tertentu,
paling rumit. Jika
$\left[ \begin{smallmatrix} f(x,y) \\ g(x,y) \end{smallmatrix} \right]$
nol, maka di sekitarnya vektor dapat menunjuk ke arah mana pun.
Selain itu, lintasan-lintasan menuju, menjauhi, atau mengitari
titik-titik ini. Jadi, jika kita mencari perilaku kualitatif jangka panjang dari sistem,
kita perlu melihat apa yang terjadi di dekat titik kritis.

Titik kritis juga kadang-kadang disebut
\emph{kesetimbangan}\index{kesetimbangan}, karena pada titik kritis terdapat apa yang disebut
\emph{solusi kesetimbangan}\index{solusi kesetimbangan}.
Jika $(x_0,y_0)$ adalah titik kritis, maka kita memiliki solusi
\begin{equation*}
x(t) = x_0, \quad y(t) = y_0 .
\end{equation*}
Dalam \examplevref{example:nlin-1b-example}, terdapat dua solusi
kesetimbangan:
\begin{equation*}
x(t) = 0, \quad y(t) = 0,
\qquad \text{dan} \qquad
x(t) = 1, \quad y(t) = 0.
\end{equation*}
Bandingkan pembahasan tentang kesetimbangan ini dengan pembahasan dalam
\sectionref{auteq:section}. Konsep yang mendasarinya
persis sama.

\subsection{Linearisasi}

Dalam \sectionref{sec:twodimaut} kita mempelajari perilaku sistem linear
homogen yang terdiri atas dua persamaan di dekat titik kritis. Untuk sistem linear
dua variabel yang diberikan oleh matriks invertibel, satu-satunya titik kritis adalah
titik asal $(0,0)$.
Mari kita manfaatkan pemahaman yang kita peroleh dalam bagian itu
untuk memahami apa yang terjadi di dekat titik kritis sistem nonlinear.

Dalam kalkulus, kita belajar menaksir suatu fungsi dengan mengambil
turunannya dan melakukan linearisasi. Kita bekerja dengan cara serupa pada sistem persamaan diferensial biasa nonlinear.
Misalkan $(x_0,y_0)$ adalah titik kritis.
Pertama, ganti variabel menjadi $(u,v)$, sehingga $(u,v)=(0,0)$ bersesuaian dengan
$(x_0,y_0)$. Artinya,
\begin{equation*}
u=x-x_0, \qquad v=y-y_0 .
\end{equation*}
Selanjutnya kita perlu mencari turunannya. Dalam kalkulus multivariabel, Anda mungkin
pernah melihat bahwa versi turunan untuk beberapa variabel adalah
\emph{\myindex{matriks Jacobian}}%
\footnote{Dinamai menurut matematikawan Jerman
\href{https://en.wikipedia.org/wiki/Carl_Gustav_Jacob_Jacobi}{Carl Gustav Jacob Jacobi}
(1804--1851).}. Matriks Jacobian dari
fungsi bernilai vektor
$\left[ \begin{smallmatrix} f(x,y) \\ g(x,y) \end{smallmatrix} \right]$
pada $(x_0,y_0)$ adalah
\begin{equation*}
\begin{bmatrix}
\frac{\partial f}{\partial x}(x_0,y_0) &
\frac{\partial f}{\partial y}(x_0,y_0) \\
\frac{\partial g}{\partial x}(x_0,y_0) &
\frac{\partial g}{\partial y}(x_0,y_0)
\end{bmatrix} .
\end{equation*}
Matriks ini memberikan hampiran linear terbaik ketika $u$ dan $v$ (dan
karena itu $x$ dan $y$) berubah.
Kita mendefinisikan \emph{\myindex{linearisasi}} persamaan
\eqref{eq:nlinautn2} sebagai sistem linear
\begin{equation*}
\begin{bmatrix} u \\ v \end{bmatrix} ' =
\begin{bmatrix}
\frac{\partial f}{\partial x}(x_0,y_0) &
\frac{\partial f}{\partial y}(x_0,y_0) \\
\frac{\partial g}{\partial x}(x_0,y_0) &
\frac{\partial g}{\partial y}(x_0,y_0)
\end{bmatrix} 
\begin{bmatrix} u \\ v \end{bmatrix} .
\end{equation*}

\begin{example} \label{example:nlin-1b-examplelin}
Mari kita gunakan kembali persamaan dalam
\exampleref{example:nlin-1b-example}:
$x' = y$, $y' = -x+x^2$. Terdapat dua titik kritis, $(0,0)$
dan $(1,0)$. Matriks Jacobian pada sembarang titik adalah
\begin{equation*}
\begin{bmatrix}
\frac{\partial f}{\partial x}(x,y) &
\frac{\partial f}{\partial y}(x,y) \\
\frac{\partial g}{\partial x}(x,y) &
\frac{\partial g}{\partial y}(x,y)
\end{bmatrix} =
\begin{bmatrix}
0 & 1 \\
-1+2x & 0
\end{bmatrix}.
\end{equation*}
Pada $(0,0)$, yaitu ketika $x=0$ dan $y=0$, matriks Jacobiannya adalah
\begin{equation*}
\begin{bmatrix}
0 & 1 \\
-1 & 0
\end{bmatrix} ,
\end{equation*}
dan variabel baru kita adalah
$u=x$ dan $v=y$. Linearisasi pada $(0,0)$ adalah
\begin{equation*}
\begin{bmatrix} u \\ v \end{bmatrix} ' =
\begin{bmatrix}
0 & 1 \\
-1 & 0
\end{bmatrix}
\begin{bmatrix} u \\ v \end{bmatrix} .
\end{equation*}

Pada titik $(1,0)$, matriks Jacobiannya
adalah
$\left[ \begin{smallmatrix}
0 & 1 \\
1 & 0
\end{smallmatrix} \right]$,
dan variabel baru kita adalah
$u=x-1$ dan $v=y$.
Linearisasi pada $(1,0)$ adalah
\begin{equation*}
\begin{bmatrix} u \\ v \end{bmatrix} ' =
\begin{bmatrix}
0 & 1 \\
1 & 0
\end{bmatrix}
\begin{bmatrix} u \\ v \end{bmatrix} .
\end{equation*}

Diagram fase dari kedua linearisasi pada
titik kritis $(0,0)$ dan $(1,0)$ diberikan dalam \figurevref{fig:nlin-1b-lin}. Perhatikan
bahwa variabelnya kini $u$ dan $v$, yang berbeda pada setiap
titik kritis. Bandingkan
\figureref{fig:nlin-1b-lin} dengan \figurevref{fig:nlin-1b}, dan perhatikan terutama
perilaku di dekat titik-titik kritis.

\begin{myfig}
\capstart
%original files nlin-1b-lin-00 nlin-1b-lin-01
\diffyincludegraphics{width=6.24in}{width=9in}{nlin-1b-lin-00-01}{%
Dua plot medan arah (kisi panah) pada bidang uv.
Di sebelah kiri, panah-panah sekadar mengitari titik asal searah jarum jam, dengan beberapa
lintasan melingkar digambar. Di sebelah kanan terdapat
gambar pelana: sepanjang garis v sama dengan u, panah-panah menjauhi
titik asal, sedangkan sepanjang v sama dengan minus u, panah-panah menuju titik asal. Garis
v sama dengan u
dan v sama dengan minus u digambar sebagai lintasan. Pada gambar yang dibagi oleh garis-garis ini,
di daerah sebelah kiri, panah-panah umumnya mengarah ke bawah, sedikit ke
kanan di bagian atas dan sedikit ke kiri di bagian bawah. Ada suatu lintasan di daerah ini
yang berawal di luar gambar pada sisi kiri,
bergerak turun dan agak ke kanan sebelum berbelok ke kiri
dan meninggalkan gambar. Situasinya serupa pada tiga daerah lainnya:
di daerah atas panah-panah kira-kira mengarah ke kanan, di daerah
kanan panah-panah kira-kira mengarah ke atas, dan di daerah bawah
panah-panah kira-kira mengarah ke kiri.}
\caption{Diagram fase beserta beberapa lintasan dari
linearisasi pada titik kritis $(0,0)$ (kiri) dan $(1,0)$ (kanan) untuk
$x' = y$, $y' = -x+x^2$. \label{fig:nlin-1b-lin}}
\end{myfig}
\end{example}

\subsection{Latihan}

\begin{exercise}
Sketsalah medan vektor bidang fase untuk:
\begin{tasks}(3)
\task $x'=x^2$, \enspace $y'=y^2$,
\task $x'=(x-y)^2$, \enspace $y'=-x$,
\task $x'=e^y$, \enspace $y'=e^x$.
\end{tasks}
\end{exercise}

\begin{samepage}
\begin{exercise}
Pasangkan sistem
\begin{tasks}[label=\arabic*)](3)
\task $x'=x^2$, \enspace $y'=y^2$,
\task $x'=xy$, \enspace $y'=1+y^2$,
\task $x'=\sin(\pi y)$, \enspace $y'=x$,
\end{tasks}
dengan medan vektor di bawah ini. Berikan alasan.
\begin{tasks}(3)
\task
\parbox[c]{1.75in}{\diffyincludegraphics{width=1.75in}{width=1.75in}{nlin-exer-xy-1py2}{%
Medan arah pada bidang xy, dengan semua panah
tampak kecil dan menunjuk langsung ke atas sepanjang sumbu x. Panah-panah di atas
sumbu x menunjuk ke atas dan menjauhi titik asal, sedangkan panah-panah di bawah
sumbu x juga menunjuk kira-kira ke atas dan menuju titik asal.}}
\task
\parbox[c]{1.75in}{\diffyincludegraphics{width=1.75in}{width=1.75in}{nlin-exer-x2-y2}{%
Medan arah pada bidang xy, dengan panah nol di titik asal,
panah pada sumbu x menunjuk ke kanan, dan panah pada sumbu y menunjuk ke atas. Di kuadran
pertama, panah-panah kira-kira menjauhi titik asal, dan di kuadran ketiga
panah-panah kira-kira menuju titik asal. Di kuadran kedua,
panah-panah mula-mula mengarah ke kanan dan perlahan berbelok ke atas ketika kita bergerak
ke kanan. Di kuadran keempat, panah-panah mula-mula mengarah ke atas di bagian bawah lalu
berbelok ke kanan ketika kita bergerak naik.}}
\task
\parbox[c]{1.75in}{\diffyincludegraphics{width=1.75in}{width=1.75in}{nlin-exer-sinpiy-x}{%
Medan arah pada bidang xy, dengan panah-panah yang tampak mengitari
dua titik pada sumbu y berlawanan arah jarum jam, satu di atas titik asal dan
satu di bawah titik asal.}}
\end{tasks}
\end{exercise}
\end{samepage}


\begin{exercise}
Carilah titik-titik kritis dan linearisasi sistem-sistem berikut.
\begin{tasks}(2)
\task $x'=x^2-y^2$, \enspace $y'=x^2+y^2-1$,
\task $x'=-y$, \enspace $y'=3x+yx^2$,
\task $x'=x^2+y$, \enspace $y'=y^2+x$.
\end{tasks}
\end{exercise}

\begin{exercise}
Untuk sistem-sistem berikut, verifikasilah bahwa $(0,0)$ merupakan titik kritis,
dan carilah linearisasinya pada $(0,0)$.
\begin{tasks}(2)
\task $x'=x+2y+x^2-y^2$, \enspace $y'=2y-x^2$
\task $x'=-y$, \enspace $y'=x-y^3$
\task* $x'=ax+by+f(x,y)$, $y'=cx+dy+g(x,y)$, dengan
$f(0,0) = 0$,
$g(0,0) = 0$, dan semua turunan parsial pertama $f$ dan $g$
juga nol pada $(0,0)$, yaitu,
$\frac{\partial f}{\partial x}(0,0) = 
\frac{\partial f}{\partial y}(0,0) = 
\frac{\partial g}{\partial x}(0,0) = 
\frac{\partial g}{\partial y}(0,0) = 0$.
\end{tasks}
\end{exercise}

\begin{exercise}
Ambil $x'=(x-y)^2$, \enspace $y'=(x+y)^2$.
\begin{tasks}
\task Carilah himpunan titik kritis.
\task Sketsalah diagram fase dan jelaskan perilaku di dekat titik(-titik)
kritis.
\task Carilah linearisasinya. Apakah linearisasi itu membantu memahami sistem?
\end{tasks}
\end{exercise}

\begin{exercise}
Ambil $x'=x^2$, \enspace $y'=x^3$.
\begin{tasks}
\task Carilah himpunan titik kritis.
\task Sketsalah diagram fase dan jelaskan perilaku di dekat titik(-titik)
kritis.
\task Carilah linearisasinya. Apakah linearisasi itu membantu memahami sistem?
\end{tasks}
\end{exercise}

\setcounter{exercise}{100}

\pagebreak[2]
\begin{exercise}
Carilah titik-titik kritis dan linearisasi sistem-sistem berikut.
\begin{tasks}(2)
\task $x'=\sin(\pi y)+(x-1)^2$, \enspace $y'=y^2-y$,
\task $x'=x+y+y^2$, \enspace $y'=x$,
\task $x'=(x-1)^2+y$, \enspace $y'=x^2+y$.
\end{tasks}
\end{exercise}
\exsol{%
a) Titik kritis $(1,0)$ dan $(1,1)$.
Pada $(1,0)$, dengan $u=x-1$, $v=y$, linearisasinya adalah
$u'=\pi v$, $v'=-v$.
Pada $(1,1)$, dengan $u=x-1$, $v=y-1$, linearisasinya adalah
$u'=-\pi v$, $v'=v$.\\
b) Titik kritis $(0,0)$ dan $(0,-1)$.
Pada $(0,0)$, dengan $u=x$, $v=y$, linearisasinya adalah
$u'=u+v$, $v'=u$.
Pada $(0,-1)$, dengan $u=x$, $v=y+1$, linearisasinya adalah
$u'=u-v$, $v'=u$.\\
c) Titik kritis $(\nicefrac{1}{2},\nicefrac{-1}{4})$. Dengan
$u=x-\nicefrac{1}{2}$, $v=y+\nicefrac{1}{4}$, linearisasinya adalah
$u'=-u+v$, $v'=u+v$.
}

\begin{exercise}
Pasangkan sistem
\begin{tasks}[label=\arabic*)](2)
\task $x'=y^2$, \enspace $y'=-x^2$,
\task $x'=y$, \enspace $y'=(x-1)(x+1)$,
\task $x'=y+x^2$, \enspace $y'=-x$,
\end{tasks}
dengan medan vektor di bawah ini. Berikan alasan.
\begin{tasks}(3)
\task
\parbox[c]{1.75in}{\diffyincludegraphics{width=1.75in}{width=1.75in}{nlin-exer-y-xm1xp1}{%
Plot medan arah pada bidang xy.
Di sebelah kiri plot, panah-panah tampak bergerak searah jarum jam mengitari suatu titik
di sebelah kiri titik asal pada sumbu x. Lalu, di sebelah kanan, panah-panah
tampak menuju suatu titik di sebelah kanan titik asal sepanjang garis diagonal
menurun, dan menjauhi titik itu sepanjang garis diagonal menaik.
Di sisi kiri plot, panah-panah kira-kira mengarah ke atas;
di bagian atas, kira-kira mengarah ke kanan; di sisi kanan, kira-kira
mengarah ke atas; dan di bagian bawah, kira-kira mengarah ke kiri.}}
\task
\parbox[c]{1.75in}{\diffyincludegraphics{width=1.75in}{width=1.75in}{nlin-exer-ypx2-mx}{%
Plot medan arah pada bidang xy.
Panah-panah tampak bergerak searah jarum jam mengitari titik asal, tetapi tidak
dengan pola yang benar-benar melingkar.
Di bagian atas, panah-panah sebagian besar mengarah ke kanan, meskipun agak mengarah ke atas
di sebelah kiri dan ke bawah di sebelah kanan, serta sedikit
lebih kecil di tengah gambar. Di sisi kiri
gambar, semua panah kira-kira menunjuk ke kiri dan ke atas,
meskipun ukurannya lebih kecil di bagian bawah gambar. Di sisi kanan
gambar, panah-panah kira-kira menunjuk ke kanan dan sedikit
ke bawah, serta lebih kecil di bagian bawah gambar. Sepanjang bagian bawah
gambar, dari kiri ke kanan, panah-panah mula-mula menunjuk ke atas dan ke kiri,
lalu perlahan berbelok ke kanan, kemudian ke bawah hingga akhirnya menunjuk
ke bawah dan ke kanan di sudut kanan bawah.}}
\task
\parbox[c]{1.75in}{\diffyincludegraphics{width=1.75in}{width=1.75in}{nlin-exer-y2-mx2}{%
Plot medan arah pada bidang xy.
Panah-panah tampak datang dari kiri atas dan kira-kira menunjuk ke bawah serta ke kanan.
Panah-panah itu terus bergerak ke arah tersebut, tetapi seakan-akan ingin
memutar mengelilingi titik asal.
Panah-panah yang berawal di atas garis y sama dengan minus x
mengikuti kurva yang mengitari bagian atas titik asal, sedangkan yang
berawal di bawah garis tersebut mengikuti kurva yang mengitari bagian bawah titik asal.}}
\end{tasks}
\end{exercise}
\exsol{%
1) adalah c), \quad 2) adalah a), \quad 3) adalah b)
}


\begin{samepage}
\begin{exercise}
Gagasan titik kritis dan linearisasi juga berlaku dalam dimensi yang lebih
tinggi. Anda cukup memperbesar matriks Jacobian dengan menambahkan lebih banyak fungsi
dan lebih banyak variabel. Untuk sistem tiga persamaan berikut,
carilah titik-titik kritis dan linearisasinya:
\begin{equation*}
x' = x + z^2, \qquad y' = z^2-y, \qquad z' = z+x^2.
\end{equation*}
\end{exercise}
\end{samepage}
\exsol{%
Titik kritisnya adalah $(0,0,0)$ dan
$(-1, 1, -1)$.
Linearisasi pada titik asal, dengan variabel $u=x$, $v=y$, $w=z$, adalah
$u' = u$, $v'=-v$, $w' = w$.
Linearisasi pada titik $(-1,1,-1)$, dengan variabel $u=x+1$,
$v=y-1$, $w=z+1$, adalah
%$x=u-1$
%$y=v+1$
%$z=w-1$
%$u' = u-1 + (w-1)^2$, $v' = (w-1)^2-v+1$, $w' = (w-1)+(u-1)^2$.
%$u' = u + w^2-2w$, $v' = w^2-2w-v$, $w' = w+u^2-2u$
$u'=u-2w$, $v'=-v-2w$, $w'=w-2u$.
}

\begin{exercise}
Setiap sistem tak otonom berdimensi dua $x'=f(x,y,t)$, $y'=g(x,y,t)$ dapat
ditulis sebagai sistem otonom berdimensi tiga (tiga persamaan). Tuliskan
sistem otonom ini menggunakan variabel $u$, $v$, $w$.
\end{exercise}
\exsol{%
$u' = f(u,v,w)$, $v'=g(u,v,w)$, $w' = 1$.
}

%%%%%%%%%%%%%%%%%%%%%%%%%%%%%%%%%%%%%%%%%%%%%%%%%%%%%%%%%%%%%%%%%%%%%%%%%%%%%%

\sectionnewpage
\section{Kestabilan dan klasifikasi titik kritis terisolasi}
\label{nlinstability:section}

%mbxINTROSUBSECTION

\sectionnotes{1,5--2 pertemuan\EPref{, \S6.1--\S6.2 dalam \cite{EP}}\BDref{,
\S9.2--\S9.3 dalam \cite{BD}}}

\subsection{Titik kritis terisolasi dan sistem hampir linear}

Suatu titik kritis disebut
\emph{terisolasi}\index{titik kritis terisolasi}
jika titik itu merupakan satu-satunya titik kritis dalam suatu
\myquote{lingkungan} kecil di sekitar titik tersebut.
Artinya, jika kita memperbesar cukup jauh, titik itulah
satu-satunya titik kritis yang terlihat.
Dalam \exampleref{example:nlin-1b-example} di atas,
kedua titik kritis terisolasi.
Sebaliknya, jika terdapat satu kurva penuh yang terdiri atas titik-titik
kritis, titik-titik itu tidak terisolasi.

Suatu sistem disebut \emph{\myindex{hampir linear}} pada titik kritis
$(x_0,y_0)$ jika titik kritis tersebut terisolasi dan matriks Jacobian pada titik itu
invertibel, atau secara ekuivalen jika sistem yang dilinearisasi memiliki titik kritis
terisolasi. Dalam kasus demikian, suku-suku nonlinear sangat kecil
dan sistem berperilaku seperti linearisasinya, setidaknya ketika kita berada dekat
dengan titik kritis.

Sebagai contoh, sistem dalam
Contoh~\ref{example:nlin-1b-example} dan \ref{example:nlin-1b-examplelin}
memiliki dua titik kritis terisolasi $(0,0)$ dan $(1,0)$, dan
hampir linear pada kedua titik kritis karena
matriks Jacobian pada kedua titik,
$\left[ \begin{smallmatrix} 0 & 1 \\ -1 & 0 \end{smallmatrix} \right]$ dan
$\left[ \begin{smallmatrix} 0 & 1 \\ 1 & 0 \end{smallmatrix} \right]$,
invertibel.

Sebaliknya, sistem $x' = x^2$, $y' = y^2$ memiliki titik kritis terisolasi
pada $(0,0)$, tetapi matriks Jacobian
\begin{equation*}
\begin{bmatrix} 2x & 0 \\ 0 & 2y \end{bmatrix}
\end{equation*}
bernilai nol ketika $(x,y) = (0,0)$. Jadi sistem ini tidak hampir
linear.
Contoh yang lebih buruk lagi adalah sistem $x' = x$, $y' = x^2$, yang tidak memiliki
titik kritis terisolasi; $x'$ dan $y'$ sama-sama nol
setiap kali $x=0$, yaitu pada seluruh sumbu $y$.

Untungnya, titik kritis paling sering
terisolasi, dan sistem hampir linear pada titik-titik kritisnya.
Jadi, jika kita mempelajari apa yang terjadi di sana, kita telah memahami sebagian besar
situasi yang muncul dalam penerapan.



\subsection{Kestabilan dan klasifikasi titik kritis terisolasi}

Setelah kita memiliki titik kritis terisolasi, sistem hampir linear pada
titik kritis tersebut, dan kita telah menghitung sistem linearisasi
yang terkait, kita dapat mengklasifikasikan perilaku
solusinya. Kurang lebih, kita menggunakan klasifikasi sistem linear
dua variabel dari \sectionref{sec:twodimaut}, dengan satu catatan
kecil.
Mari kita daftar perilakunya berdasarkan nilai eigen
matriks Jacobian pada titik kritis dalam \tablevref{pln:behtab2}.
Tabel ini sangat mirip dengan \tablevref{pln:behtab},
kecuali tidak adanya titik \myquote{pusat}.
%There is also a new column
%that we will discuss.
Kita akan membahas pusat kemudian, karena kasus itu lebih rumit.

\begin{table}[h!t]
\mybeginframe
\capstart
\begin{center}
\begin{tabular}{@{}lll@{}}
\toprule
Nilai eigen matriks Jacobian & Perilaku & Kestabilan \\
\midrule
real dan keduanya positif & sumber / simpul tak stabil & tak stabil \\
real dan keduanya negatif & serapan / simpul stabil & stabil asimtotik \\
real dan berlainan tanda & pelana & tak stabil \\
kompleks dengan bagian real positif & sumber spiral & tak stabil \\
kompleks dengan bagian real negatif & serapan spiral & stabil asimtotik \\
\bottomrule
\end{tabular}
\end{center}
\caption{Perilaku sistem hampir linear di dekat titik kritis
terisolasi. \label{pln:behtab2}}
\myendframe
\end{table}

Pada kolom ketiga,
kita menandai titik sebagai \emph{stabil asimtotik} atau \emph{tak stabil}. Secara formal,
\emph{\myindex{titik kritis stabil}} $(x_0,y_0)$ adalah titik yang memenuhi hal berikut:
untuk setiap jarak kecil $\epsilon$ dari $(x_0,y_0)$, dan setiap kondisi awal
dalam radius yang mungkin lebih kecil di sekitar $(x_0,y_0)$, lintasan
sistem tidak pernah menjauh dari $(x_0,y_0)$ melebihi $\epsilon$.
\emph{\myindex{Titik kritis tak stabil}} adalah titik yang tidak stabil.
Secara informal, suatu titik stabil jika, ketika kita mulai dekat dengan titik kritis dan
mengikuti suatu lintasan, kita menuju, atau setidaknya tidak menjauh
dari,
titik kritis tersebut.

Titik kritis stabil $(x_0,y_0)$ disebut \emph{\myindex{stabil asimtotik}} jika,
untuk setiap kondisi awal yang cukup dekat dengan $(x_0,y_0)$ dan setiap
solusi $\bigl( x(t), y(t) \bigr)$ yang memenuhi kondisi tersebut,
\begin{equation*}
\lim_{t \to \infty} \bigl( x(t), y(t) \bigr) = (x_0,y_0) .
\end{equation*}
Artinya, titik kritis stabil asimtotik
jika setiap lintasan dengan kondisi awal yang cukup dekat
menuju titik kritis $(x_0,y_0)$.

\begin{example} \label{example:nlin-xplusy}
Tinjau
$x'=-y-x^2$,
$y'=-x+y^2$.
Lihat \figurevref{fig:nlin-ex813-new} untuk diagram fasenya.
Mari kita cari titik-titik kritisnya. Inilah titik-titik tempat
$-y-x^2 = 0$ dan $-x+y^2=0$. Persamaan pertama berarti $y = -x^2$, sehingga
$y^2 = x^4$. Dengan menyubstitusikannya ke persamaan kedua, kita memperoleh
$-x+x^4 = 0$. Dengan memfaktorkan, kita memperoleh $x(x^3-1)=0$. Karena kita hanya mencari
solusi real, kita mendapatkan $x=0$ atau $x=1$. Dengan mencari $y$ yang
bersesuaian menggunakan $y = -x^2$, kita memperoleh dua titik kritis, yaitu $(0,0)$
dan $(1,-1)$. Jelas bahwa titik-titik kritis ini terisolasi.

\begin{myfig}
\capstart
\diffyincludegraphics{width=3in}{width=4.5in}{nlin-ex813-new}{%
Plot medan arah pada bidang xy untuk x antara minus 2 dan 2
serta y antara minus 2 dan 2. Terdapat dua titik kritis, satu di
titik asal dan satu di (1, minus 1). Titik asal berupa pelana, dengan panah-panah
yang datang ke arahnya sepanjang kurva diagonal menaik dan panah-panah yang keluar
sepanjang kurva diagonal menurun. Titik (1, minus 1) adalah serapan, dengan
panah-panah yang menunjuk ke arahnya. Beberapa lintasan diperlihatkan; lintasan yang
berawal di sisi kanan untuk y kira-kira di bawah
y sama dengan 1 dan
di sisi bawah kira-kira di kanan x sama dengan minus 1 bergerak masuk ke serapan. Lintasan
lain berbelok menuju dan meninggalkan gambar pada sisi kiri atas
plot.}
\caption{Potret fase beserta beberapa contoh lintasan dari
$x'=-y-x^2$, $y'=-x+y^2$. \label{fig:nlin-ex813-new}}
\end{myfig}

Matriks Jacobiannya adalah
\begin{equation*}
\begin{bmatrix}
-2x & -1 \\
-1 & 2y
\end{bmatrix} .
\end{equation*}
Pada titik $(0,0)$, kita memperoleh matriks
$\left[ \begin{smallmatrix} 0 & -1 \\ -1 & 0 \end{smallmatrix} \right]$
yang nilai eigennya $1$ dan $-1$. Karena matriksnya invertibel, sistem hampir linear
pada $(0,0)$. Karena nilai eigennya real
dan berlainan tanda, kita memperoleh titik pelana, yaitu titik kesetimbangan
yang tak stabil.

Pada titik $(1,-1)$, kita memperoleh matriks
$\left[ \begin{smallmatrix} -2 & -1 \\ -1 & -2 \end{smallmatrix} \right]$
yang
nilai eigennya $-1$ dan $-3$.
Matriksnya invertibel, sehingga sistem hampir linear pada $(1,-1)$.
Karena kedua nilai eigen real dan negatif, titik
kritis tersebut merupakan serapan, dan karena itu merupakan titik kesetimbangan yang stabil asimtotik.
Artinya, jika kita mulai dari sembarang titik $(x_i,y_i)$ yang dekat dengan $(1,-1)$ sebagai
kondisi awal dan memplot suatu lintasan, lintasan itu mendekati $(1,-1)$.
Dengan kata lain,
\begin{equation*}
\lim_{t \to \infty} \bigl( x(t), y(t) \bigr) = (1,-1) .
\end{equation*}
Seperti yang dapat
Anda lihat dari diagram, perilaku ini bahkan berlaku untuk beberapa
titik awal yang cukup jauh dari $(1,-1)$, tetapi tentu tidak berlaku untuk semua
titik awal.
\end{example}

\begin{example} \label{example:nlin-withexp}
Tinjau
$x'=y+y^2e^x$,
$y'=x$. Pertama, kita mencari titik-titik kritis. Inilah titik tempat
$y+y^2e^x = 0$ dan $x=0$. Dengan menyederhanakan, kita memperoleh $0=y+y^2 = y(y+1)$. Jadi,
titik kritisnya adalah $(0,0)$ dan $(0,-1)$, sehingga keduanya terisolasi. Mari kita
hitung matriks Jacobiannya:
\begin{equation*}
\begin{bmatrix}
y^2e^x & 1+2ye^x \\
1 & 0
\end{bmatrix}.
\end{equation*}

Pada titik $(0,0)$ kita memperoleh matriks
$\left[ \begin{smallmatrix} 0 & 1 \\ 1 & 0 \end{smallmatrix} \right]$, sehingga
kedua nilai eigennya adalah $1$ dan $-1$. Karena matriksnya invertibel, sistem hampir linear
pada $(0,0)$. Dan, karena nilai eigennya real
serta berlainan tanda, kita memperoleh titik pelana, yaitu titik kesetimbangan
yang tak stabil.

Pada titik $(0,-1)$ kita memperoleh matriks
$\left[ \begin{smallmatrix} 1 & -1 \\ 1 & 0 \end{smallmatrix} \right]$ yang
nilai eigennya adalah $\frac{1}{2} \pm i \frac{\sqrt{3}}{2}$.
Matriksnya invertibel, sehingga sistem hampir linear pada $(0,-1)$.
Karena kita memiliki nilai eigen kompleks dengan bagian real positif, titik
kritis itu merupakan sumber spiral, dan karena itu merupakan titik kesetimbangan yang tak stabil.

\begin{myfig}
\capstart
\diffyincludegraphics{width=3in}{width=4.5in}{nlin-ex813}{%
Plot medan arah pada bidang xy untuk x antara minus 2 dan 2
serta y antara minus 2 dan 2. Titik asal adalah titik kritis berupa pelana,
dengan panah-panah menuju titik asal sepanjang kurva miring menurun
dan panah-panah keluar sepanjang kurva miring menaik. Titik
(0, minus 1) merupakan sumber spiral, dengan panah-panah yang melilit keluar dari
titik kritis berlawanan arah jarum jam. Beberapa lintasan diperlihatkan, terutama dua
lintasan yang berpilin keluar dari titik kritis (0, minus 1).}
\caption{Potret fase beserta beberapa contoh lintasan dari
$x'=y+y^2e^x$, $y'=x$. \label{fig:nlin-ex813}}
\end{myfig}

Lihat \figurevref{fig:nlin-ex813} untuk diagram fasenya. Perhatikan kedua
titik kritis serta perilaku panah-panah dalam medan vektor di sekitar
titik-titik tersebut.
\end{example}

\subsection{Masalah pada pusat}

Ingatlah bahwa pusat pada sistem linear berarti lintasan-lintasan
bergerak dalam orbit eliptis tertutup
dengan suatu arah mengitari titik kritis. Titik
kritis seperti itu kita sebut \emph{\myindex{pusat}} atau
\emph{\myindex{pusat stabil}}. Titik ini bukan titik kritis yang stabil
asimtotik, karena lintasan tidak pernah mendekati titik kritis,
tetapi setidaknya jika Anda mulai cukup dekat dengan titik kritis,
Anda akan tetap dekat dengannya. Contoh paling sederhana dari
perilaku semacam itu adalah sistem linear dengan pusat. Contoh
lain adalah titik kritis $(0,0)$ dalam
\examplevref{example:nlin-1b-example}.

Masalah pada pusat dalam sistem nonlinear terletak pada kenyataan bahwa
gerak suatu lintasan menuju atau menjauhi titik kritis ditentukan oleh tanda
bagian real nilai eigen matriks Jacobian, sedangkan matriks Jacobian
dalam sistem nonlinear berubah dari titik ke titik. Karena bagian
real ini nol pada titik kritis itu sendiri, di dekatnya nilainya dapat bertanda apa pun,
yang berarti lintasan dapat ditarik menuju atau menjauhi titik
kritis.

\begin{example}
Contoh perilaku bermasalah semacam itu adalah sistem
$x'=y, y' = -x+y^3$. Satu-satunya titik kritis
adalah titik asal $(0,0)$. Matriks Jacobiannya adalah
\begin{equation*}
\begin{bmatrix}
0 & 1 \\
-1 & 3 y^2 \\
\end{bmatrix} .
\end{equation*}
Pada
$(0,0)$, matriks Jacobiannya adalah
$\left[ \begin{smallmatrix}
0 & 1 \\
-1 & 0 \\
\end{smallmatrix} \right]$, yang memiliki nilai eigen $\pm i$. Jadi,
linearisasinya memiliki pusat.

Dengan menggunakan persamaan kuadrat, nilai eigen
matriks Jacobian pada sembarang titik $(x,y)$ adalah
\begin{equation*}
\lambda = 
\frac{3}{2}y^2 \pm
i
\frac{\sqrt{4-9y^4}}{2} .
\end{equation*}
Pada setiap titik dengan $y \neq 0$ (jadi, pada hampir semua titik di dekat titik asal), nilai eigennya memiliki bagian real positif ($y^2$ tidak pernah
negatif). Bagian real positif ini
mendorong lintasan menjauhi titik asal. Contoh lintasan untuk suatu
kondisi awal di dekat titik asal diberikan dalam
\figurevref{fig:nlin-unstable-center}.
\begin{myfig}
\capstart
\diffyincludegraphics{width=3in}{width=4.5in}{nlin-unstable-centerfig}{%
Plot medan arah pada bidang xy.
Panah-panah di sekitar titik asal tampak mengitari titik asal
berlawanan arah jarum jam. Diperlihatkan suatu lintasan yang sangat rapat melilit,
yang melilit keluar dari titik kritis dan akhirnya meninggalkan gambar (bahkan
berbelok ke arah yang berbeda dari arah lilitannya).}
\caption{Titik kritis tak stabil (sumber spiral) pada titik asal
untuk $x'=y, y' = -x+y^3$, meskipun linearisasinya memiliki pusat. \label{fig:nlin-unstable-center}}
\end{myfig}
\end{example}

Pesan dari contoh tersebut adalah bahwa analisis lebih lanjut diperlukan ketika
linearisasi memiliki pusat. Analisisnya secara umum akan lebih
rumit daripada contoh di atas dan lebih mungkin memerlukan
pertimbangan kasus demi kasus. Kerumitan seperti ini seharusnya tidak
mengejutkan Anda. Pada tahap ini dalam perjalanan matematika Anda, Anda telah
menemui banyak keadaan ketika suatu uji sederhana tidak memberikan kesimpulan,
sehingga diperlukan analisis yang lebih cermat dan mungkin bersifat ad hoc.
Ingatlah, misalnya, ketika
uji turunan kedua untuk maksimum atau minimum
tidak memberikan kesimpulan.

\subsection{Persamaan konservatif}

Persamaan berbentuk
\begin{equation*}
x'' + f(x) = 0,
\end{equation*}
dengan $f(x)$ fungsi sembarang, disebut
\emph{\myindex{persamaan konservatif}}. Sebagai contoh, persamaan bandul
adalah persamaan konservatif. Persamaan-persamaan ini konservatif karena tidak ada
gesekan dalam sistem sehingga energi dalam sistem \myquote{kekal.}
Mari kita tuliskan persamaan ini sebagai
sistem PDB nonlinear:
\begin{equation*}
x' = y, \qquad y' = -f(x) .
\end{equation*}
Jenis persamaan ini memiliki
keuntungan bahwa kita dapat mencari lintasannya dengan mudah.

Caranya adalah mula-mula memandang $y$ sebagai fungsi $x$ untuk sementara. Lalu,
gunakan aturan rantai
\begin{equation*}
x'' = y' = \frac{dy}{dx} x' = y \frac{dy}{dx} ,
\end{equation*}
dengan tanda prima menyatakan turunan terhadap $t$.
Kita memperoleh $y \frac{dy}{dx} + f(x) = 0$. Kita mengintegralkan terhadap
$x$ untuk mendapatkan
$\int y \frac{dy}{dx} \,dx + \int f(x)\, dx = C$. Dengan kata lain,
\begin{equation*}
\frac{1}{2} y^2  + \int f(x)\, dx = C .
\end{equation*}
Kita memperoleh persamaan implisit untuk lintasan, dengan nilai $C$ yang berbeda
memberikan lintasan yang berbeda. Nilai
$C$ tetap kekal pada setiap lintasan. Ungkapan ini
kadang-kadang disebut \emph{\myindex{Hamiltonian}} atau energi
sistem.
Jika Anda melihat kembali \sectionref{exact:section}, Anda akan memperhatikan
bahwa $y\frac{dy}{dx} + f(x) = 0$ adalah persamaan eksak, dan
kita baru saja menemukan fungsi potensial.

\begin{example}
Mari kita cari lintasan untuk persamaan $x'' + x-x^2 = 0$,
yakni persamaan dari
\examplevref{example:nlin-1b-example}. Sistem orde pertama
yang bersesuaian adalah
\begin{equation*}
x' = y , \qquad y' = -x+x^2 .
\end{equation*}
Lintasan memenuhi
\begin{equation*}
\frac{1}{2} y^2  + \frac{1}{2} x^2 - \frac{1}{3} x^3  = C .
\end{equation*}
Kita selesaikan untuk $y$
\begin{equation*}
y = \pm \sqrt{-x^2 + \frac{2}{3} x^3  + 2C} .
\end{equation*}

Dengan memplot grafik-grafik ini, kita mendapatkan tepat lintasan-lintasan dalam
\figurevref{fig:nlin-1b}. Secara khusus, kita melihat bahwa di dekat titik asal
lintasannya berupa \emph{\myindex{kurva tertutup}}: lintasan itu terus
mengitari titik asal, tanpa pernah berpilin masuk atau keluar. Dengan demikian, kita menemukan cara
untuk memverifikasi bahwa titik kritis pada $(0,0)$ adalah pusat stabil.
Titik kritis pada $(1,0)$ adalah pelana, seperti yang telah kita perhatikan.
Contoh ini merupakan ciri khas persamaan konservatif.
\end{example}

Tinjau persamaan konservatif sembarang
$x'' + f(x) = 0$.
Semua titik kritis terjadi ketika $y=0$ (pada
sumbu $x$), yaitu ketika $x' = 0$. Titik kritisnya adalah
titik-titik pada sumbu $x$ tempat $f(x) = 0$.
Lintasannya diberikan oleh
\begin{equation*}
y = \pm \sqrt{ - 2 \int f(x)\, dx + 2C} .
\end{equation*}
Jadi, semua lintasan simetris terhadap sumbu $x$. Secara khusus,
tidak mungkin ada sumber atau serapan spiral.
Matriks Jacobiannya adalah
\begin{equation*}
\begin{bmatrix}
0 & 1 \\
-f'(x) & 0
\end{bmatrix} .
\end{equation*}
Titik kritisnya hampir linear jika $f'(x) \neq 0$ pada titik
kritis. Misalkan $J$ menyatakan matriks Jacobian.
Nilai eigen $J$ adalah solusi dari
\begin{equation*}
0 = \det(J - \lambda I) = \lambda^2 + f'(x) .
\end{equation*}
Karena itu, $\lambda = \pm \sqrt{-f'(x)}$. Dengan kata lain, kita memperoleh
nilai eigen real yang berlainan tanda (jika $f'(x) < 0$),
atau nilai eigen imajiner murni (jika $f'(x) > 0$).
Hanya ada dua kemungkinan untuk titik kritis, yaitu \emph{titik pelana
tak stabil}, atau \emph{pusat stabil}.
Tidak pernah ada serapan atau sumber.

\subsection{Latihan}

\begin{exercise}
Untuk sistem di bawah ini, cari dan klasifikasikan titik-titik kritisnya.
Nyatakan pula
apakah kesetimbangannya stabil, stabil asimtotik, atau tak stabil.
\begin{tasks}(2)
\task $x'=-x+3x^2, y'=-y$
\task $x'=x^2+y^2-1$, $y'=x$
\task $x'=ye^x$, $y'=y-x+y^2$
\end{tasks}
\end{exercise}

\begin{exercise}
\pagebreak[2]
Carilah persamaan implisit lintasan untuk sistem konservatif
berikut. Selanjutnya, carilah titik kritisnya (jika ada) dan klasifikasikan.
\begin{tasks}(2)
\task $x''+ x+x^3 = 0$
\task $\theta''+\sin \theta = 0$
\task $z''+ (z-1)(z+1) = 0$
\task $x''+ x^2+1 = 0$
\end{tasks}
\end{exercise}

\begin{exercise}
Cari dan klasifikasikan titik(-titik) kritis dari $x' = -x^2$, $y' = -y^2$.
\end{exercise}

\begin{samepage}
\begin{exercise}
Misalkan $x'=-xy$, $y'=x^2-1-y$.
\begin{tasks}
\task
Tunjukkan bahwa terdapat dua serapan spiral di
$(-1,0)$ dan $(1,0)$.
\task
Untuk setiap titik awal berbentuk $(0,y_0)$, carilah lintasannya.
\task
Dapatkah suatu lintasan yang berawal pada $(x_0,y_0)$ dengan $x_0 > 0$ berpilin menuju
titik kritis $(-1,0)$? Mengapa atau mengapa tidak?
\end{tasks}
\end{exercise}
\end{samepage}

\begin{exercise} \label{exercise:increasing}
Dalam contoh $x'=y$, $y'=y^3-x$,
tunjukkan bahwa untuk setiap lintasan, jarak
dari titik asal merupakan fungsi yang meningkat.
Simpulkan
bahwa pada titik asal sistem berperilaku seperti sumber spiral.
Petunjuk: Tinjau $f(t) =
{\bigl(x(t)\bigr)}^2 + 
{\bigl(y(t)\bigr)}^2$ dan tunjukkan bahwa turunannya positif.
\end{exercise}


\begin{exercise}
Misalkan $f$ selalu positif.
Carilah lintasan dari $x''+f(x') = 0$.
Adakah titik kritis?
\end{exercise}

\begin{exercise}
Misalkan $x' = f(x,y)$, $y' = g(x,y)$. Misalkan $g(x,y) > 1$ untuk
semua $x$ dan $y$. Adakah titik kritis? Apa yang dapat kita katakan tentang
lintasan ketika $t$ menuju tak hingga?
\end{exercise}

\setcounter{exercise}{100}

\begin{exercise}
Untuk sistem-sistem di bawah ini, cari dan klasifikasikan titik-titik kritisnya.
\begin{tasks}(3)
\task $x'=-x+x^2$, $y'=y$
\task $x'=y-y^2-x$, $y'=-x$
\task $x'=xy$, $y'=x+y-1$
\end{tasks}
\end{exercise}
\exsol{%
a) $(0,0)$: pelana (tak stabil), $(1,0)$: sumber (tak stabil), \qquad
b) $(0,0)$: serapan spiral (stabil asimtotik), $(0,1)$: pelana (tak stabil), \qquad
c) $(1,0)$: pelana (tak stabil), $(0,1)$: sumber (tak stabil)
}

\begin{exercise}
Carilah persamaan implisit lintasan untuk sistem konservatif
berikut. Selanjutnya, carilah titik kritisnya (jika ada) dan klasifikasikan.
\begin{tasks}(3)
\task $x''+ x^2 = 4$
\task $x''+ e^x = 0$
\task $x''+ (x+1)e^x = 0$
\end{tasks}
\end{exercise}
\exsol{%
a) $\frac{1}{2}y^2 + \frac{1}{3}x^3 -4x = C$, titik kritis:
$(-2,0)$, pelana tak stabil, dan $(2,0)$, pusat stabil. \quad
b) $\frac{1}{2}y^2 + e^x = C$, tidak ada titik kritis. \quad
c) $\frac{1}{2}y^2 + xe^x = C$, titik kritis pada $(-1,0)$ adalah pusat stabil.
}

\begin{exercise}
Sistem konservatif $x''+x^3 = 0$ tidak hampir linear. Meskipun demikian,
klasifikasikan titik(-titik) kritisnya.
\end{exercise}
\exsol{%
Titik kritis pada $(0,0)$.
Lintasannya adalah $y = \pm \sqrt{2C-(\nicefrac{1}{2})x^4}$; untuk $C > 0$, lintasan ini menghasilkan kurva
tertutup di sekitar titik asal, sehingga titik kritis tersebut merupakan pusat stabil.
}

\begin{exercise}
Turunkan klasifikasi yang analog untuk titik kritis pada persamaan berdimensi satu,
seperti $x'= f(x)$, berdasarkan turunannya. Titik $x_0$ kritis ketika $f(x_0) = 0$ dan
hampir linear jika sebagai tambahan $f'(x_0) \neq 0$. Tentukan apakah titik kritis itu stabil atau tak stabil
berdasarkan tanda $f'(x_0)$. Jelaskan. Petunjuk: lihat \sectionref{auteq:section}.
\end{exercise}
\exsol{%
Titik kritis $x_0$ stabil jika $f'(x_0) < 0$ dan tak stabil jika $f'(x_0)
> 0$.
}

%%%%%%%%%%%%%%%%%%%%%%%%%%%%%%%%%%%%%%%%%%%%%%%%%%%%%%%%%%%%%%%%%%%%%%%%%%%%%%


\sectionnewpage
\section{Penerapan sistem nonlinear}
\label{nlinapps:section}

%mbxINTROSUBSECTION

\sectionnotes{2 pertemuan\EPref{, \S6.3--\S6.4 dalam \cite{EP}}\BDref{, \S9.3,
\S9.5 dalam \cite{BD}}}

Dalam bagian ini kita mempelajari dua contoh sistem nonlinear yang sangat
umum. Pertama, kita meninjau persamaan bandul nonlinear. Sebelumnya kita telah
melihat linearisasi persamaan bandul, tetapi kita mencatat bahwa linearisasi itu
hanya berlaku untuk sudut kecil dan waktu singkat. Sekarang kita mencari tahu
apa yang terjadi untuk sudut besar. Berikutnya, kita meninjau persamaan
pemangsa--mangsa, yang memiliki berbagai penerapan dalam persoalan pemodelan
di bidang biologi, kimia, ekonomi, dan bidang lainnya.

\subsection{Bandul}

Contoh pertama yang kita pelajari adalah persamaan bandul
$\theta''+\frac{g}{L} \sin \theta = 0$. Di sini, $\theta$ adalah simpangan
sudut, $g$ adalah percepatan gravitasi, dan $L$ adalah panjang bandul. Dalam
persamaan ini kita mengabaikan gesekan, jadi yang kita bicarakan adalah bandul
ideal.

\begin{mywrapfigsimp}{1.45in}{1.75in}
\noindent
\inputpdft{mv-pend}{%
%note that this diagram appears more than once in the book
Diagram sebuah bandul. Sebuah bandul tampak terikat pada satu titik di bagian
atas dan bergerak sepanjang busur lingkaran. Bagian bawah bandul diberi label
m, dan sudut antara bandul dan garis vertikal diberi label theta.
Panjang bandul diberi label L.}
\end{mywrapfigsimp}
Persamaan ini adalah persamaan konservatif, sehingga kita dapat menggunakan
analisis persamaan konservatif dari bagian sebelumnya.
Kita mengubah persamaan tersebut menjadi sistem dua dimensi
dalam variabel $(\theta,\omega)$ dengan memperkenalkan variabel baru
$\omega = \theta'$:
\begin{equation*}
\begin{bmatrix}
\theta \\ \omega
\end{bmatrix} '
=
\begin{bmatrix}
\omega \\
- \frac{g}{L} \sin \theta
\end{bmatrix} .
\end{equation*}
Titik kritis sistem ini terjadi ketika $\omega = 0$ dan $-\frac{g}{L}
\sin \theta = 0$, atau dengan kata lain ketika $\sin \theta = 0$. Jadi titik
kritisnya terjadi ketika $\omega = 0$ dan $\theta$ merupakan kelipatan
$\pi$. Artinya, titik-titik tersebut adalah $\ldots (-2\pi,0), (-\pi,0),
(0,0), (\pi,0), (2\pi,0) \ldots$. Walaupun terdapat tak hingga banyak titik
kritis, semuanya terisolasi. Untuk persamaan konservatif, terdapat dua jenis
titik kritis: pusat stabil atau titik pelana.

Kita menghitung matriks Jacobian:
\begin{equation*}
\begin{bmatrix}
\frac{\partial}{\partial \theta} \Bigl( \omega \Bigr) & 
\frac{\partial}{\partial \omega} \Bigl( \omega \Bigr) \\
\frac{\partial}{\partial \theta} \Bigl( - \frac{g}{L} \sin \theta \Bigr) & 
\frac{\partial}{\partial \omega} \Bigl( - \frac{g}{L} \sin \theta \Bigr)
\end{bmatrix}
=
\begin{bmatrix}
0 & 1 \\
- \frac{g}{L} \cos \theta & 0
\end{bmatrix} .
\end{equation*}
Nilai eigen matriks ini adalah
$\lambda = \pm \sqrt{-\frac{g}{L}\cos \theta}$.
Nilai eigennya bernilai real ketika $\cos \theta < 0$. Hal ini terjadi pada
kelipatan ganjil $\pi$. Nilai eigennya murni imajiner ketika
$\cos \theta > 0$. Hal ini terjadi pada kelipatan genap $\pi$. Oleh karena
itu, sistem memiliki pusat stabil di titik-titik $\ldots (-2\pi,0), (0,0),
(2\pi,0) \ldots$, dan memiliki pelana tak stabil di titik-titik
$\ldots (-3\pi,0), (-\pi,0), (\pi,0), (3\pi,0) \ldots$. Perhatikan diagram
fase pada \figurevref{fig:nlin-pend-phasediag}, tempat kita mengambil
$\frac{g}{L} = 1$ demi kesederhanaan.

\begin{myfig}
\capstart
\diffyincludegraphics{width=3in}{width=4.5in}{nlin-pend-phasediag}{%
Plot medan arah pada bidang theta omega
(dengan sudut theta sebagai arah horizontal)
untuk theta antara minus 7 dan 7 serta omega antara minus 3 dan 3.
Tampak 5 titik kritis.
Tampak 3 pusat, satu di titik asal, satu di (minus 2 pi,0),
dan satu di (2 pi,0). Anak-anak panah berputar searah jarum jam mengelilingi
pusat-pusat tersebut. Tampak 2 pelana, satu di (minus pi,0) dan satu di
(pi,0). Beberapa lintasan digambar. Lintasan di sekitar pusat hanya
mengelilinginya dalam bentuk elips. Lintasan yang bermula tinggi di sisi
kiri bergerak dengan gerak berombak seolah-olah ditolak oleh pusat-pusat
sepanjang bagian atas plot dan keluar di sisi kanan plot.
Lintasan di bagian bawah melakukan hal yang sama dengan arah terbalik.
Ada pula lintasan yang mengelilingi pusat-pusat dan bertemu
pelana secara diagonal.}
\caption{Diagram bidang fase dan beberapa lintasan persamaan
bandul nonlinear. \label{fig:nlin-pend-phasediag}}
\end{myfig}

Dalam persamaan yang dilinearisasi, kita hanya memiliki satu titik kritis,
yaitu pusat di $(0,0)$. Sekarang kita melihat dengan lebih jelas apa yang
dimaksud ketika kita mengatakan bahwa linearisasi baik untuk sudut kecil.
Sumbu horizontal adalah sudut simpangan. Sumbu vertikal adalah kecepatan
sudut bandul. Misalkan kita mulai dari $\theta = 0$ (tanpa simpangan), dan
memulai dengan kecepatan sudut $\omega$ yang kecil. Lintasan kemudian terus
mengelilingi titik kritis $(0,0)$ sepanjang suatu lingkaran hampiran. Hal ini bersesuaian
dengan ayunan pendek bandul ke depan dan ke belakang. Ketika $\theta$ tetap
kecil, lintasan benar-benar tampak seperti lingkaran sehingga sangat dekat
dengan linearisasi kita.

Ketika kita memberi bandul dorongan yang cukup besar, bandul melewati titik
teratas dan terus berputar pada porosnya. Perilaku ini bersesuaian dengan
kurva-kurva berombak yang tidak memotong sumbu horizontal pada diagram fase.
Misalkan kita meninjau kurva-kurva atas, ketika kecepatan sudut $\omega$
besar dan positif. Saat itu bandul terus berputar mengelilingi porosnya.
Kecepatannya besar ketika bandul berada dekat titik terbawah, dan paling kecil
ketika bandul berada dekat titik teratas lintasannya.

Pada setiap titik kritis terdapat solusi kesetimbangan. Pertimbangkan solusi
$\theta = 0$; bandul tidak bergerak dan menggantung lurus ke bawah. Posisi
ini stabil bagi bandul, sehingga merupakan kesetimbangan \emph{stabil}.

Jenis solusi kesetimbangan lainnya berada pada titik tak stabil, misalnya
$\theta = \pi$. Di sini bandul terbalik. Memang, bandul dapat diseimbangkan
dalam posisi ini dan akan tetap di sana, tetapi ini merupakan kesetimbangan
\emph{tak stabil}. Dorongan sekecil apa pun akan membuat bandul mulai berayun
dengan hebat.

Lihat \figurevref{fig:nlin-pend} untuk diagramnya. Gambar pertama adalah
kesetimbangan stabil $\theta = 0$. Gambar kedua bersesuaian dengan
\myquote{lingkaran-lingkaran hampiran} pada diagram fase di sekitar
$\theta =0$ ketika kecepatan sudut kecil. Gambar ketiga adalah kesetimbangan
tak stabil $\theta = \pi$. Gambar terakhir bersesuaian dengan garis-garis
berombak untuk kecepatan sudut besar.

\begin{myfig}
\capstart
\inputpdft{nlin-pend}{%
Empat diagram kemungkinan gerak bandul.
Dari kiri ke kanan, gambar pertama menunjukkan bandul diam yang menggantung
ke bawah dengan label 'theta sama dengan 0'.
Gambar kedua, berlabel 'Kecepatan sudut kecil', menunjukkan bandul berayun
dengan anak panah yang menandai gerak bolak-balik pada busur lingkaran.
Gambar ketiga, berlabel 'theta sama dengan pi', menunjukkan bandul diam yang
seimbang tepat ke atas. Gambar terakhir, berlabel 'Kecepatan sudut besar',
menunjukkan bandul yang berputar penuh mengelilingi porosnya.}
\caption{Berbagai kemungkinan gerak bandul. \label{fig:nlin-pend}}
\end{myfig}

Besaran
\begin{equation*}
\frac{1}{2} \omega^2  - \frac{g}{L} \cos \theta 
\avoidbreak
\end{equation*}
tetap konstan pada setiap solusi. Besaran ini adalah energi atau Hamiltonian
sistem.

Karena persamaan kita konservatif, maka (sebagai latihan)
lintasannya diberikan oleh
\begin{equation*}
\omega = \pm \sqrt{ \frac{2g}{L} \cos \theta + C} ,
\end{equation*}
untuk berbagai nilai $C$.
Mari kita tinjau kondisi awal $(\theta_0,0)$, yaitu kita membawa bandul ke
sudut $\theta_0$, lalu melepaskannya begitu saja (kecepatan sudut awal 0).
Kita substitusikan kondisi awal ke persamaan di atas dan menyelesaikannya
untuk $C$, sehingga diperoleh
\begin{equation*}
C = - \frac{2g}{L} \cos \theta_0 .
\end{equation*}
Jadi, rumus lintasannya adalah
\begin{equation*}
\omega = \pm \sqrt{ \frac{2g}{L}} \sqrt{ \cos \theta - \cos \theta_0 } .
\end{equation*}

Mari kita tentukan periodenya, yaitu waktu yang diperlukan bandul untuk
berayun ke depan dan ke belakang.
Kita perhatikan bahwa lintasan di sekitar titik asal pada bidang fase
simetris terhadap sumbu $\theta$ maupun sumbu $\omega$. Artinya, ditinjau
dari $\theta$, waktu yang diperlukan dari $\theta_0$ ke $-\theta_0$ sama
dengan waktu dari $-\theta_0$ kembali ke $\theta_0$. Selain itu, waktu yang
diperlukan dari $-\theta_0$ ke $0$ sama dengan waktu dari $0$ ke
$\theta_0$. Oleh karena itu, mari kita cari waktu yang diperlukan bandul
untuk bergerak dari sudut 0 ke sudut $\theta_0$, yaitu seperempat osilasi
penuh, lalu kita kalikan dengan 4.

Kita menentukan waktu ini dengan mencari
$\frac{dt}{d\theta}$ dan mengintegralkannya dari $0$ sampai $\theta_0$.
Periode adalah empat kali integral ini. Mari kita tetap berada di daerah
tempat $\omega$ positif. Karena $\omega = \frac{d\theta}{dt}$, dengan
membaliknya kita memperoleh
\begin{equation*}
\frac{dt}{d\theta} = \sqrt{\frac{L}{2g}} \frac{1}{\sqrt{\cos \theta - \cos \theta_0 }} .
\end{equation*}
Oleh karena itu, periode $T$ diberikan oleh
\begin{equation*}
T = 4 \sqrt{\frac{L}{2g}} \int_0^{\theta_0} \frac{1}{\sqrt{\cos \theta -
\cos \theta_0 }}\, d\theta .
\end{equation*}
Integral tersebut merupakan integral tak wajar, dan secara umum tidak dapat
kita evaluasi secara simbolik. Jika ingin menghitung suatu nilai $T$
tertentu, kita harus menggunakan hampiran numerik.

Ingat dari \sectionref{sec:mv}, persamaan yang dilinearisasi
$\theta''+\frac{g}{L}\theta = 0$ memiliki periode
\begin{equation*}
T_{\text{linear}} = 2\pi \sqrt{\frac{L}{g}} .
\end{equation*}
Kita memplot $T$, $T_{\text{linear}}$, dan galat relatif
$\frac{T-T_{\text{linear}}}{T}$ pada \figurevref{fig:TvsT0}. Galat relatif
menyatakan, dalam persentase, seberapa jauh hampiran kita dari periode
sebenarnya. Perhatikan bahwa $T_{\text{linear}}$ hanyalah sebuah konstanta.
Nilainya tidak berubah bersama sudut awal $\theta_0$. Periode sebenarnya $T$
semakin besar ketika $\theta_0$ semakin besar.
Perhatikan bahwa galat relatif kecil ketika $\theta_0$ kecil. Galatnya masih
hanya $15\%$ ketika $\theta_0 = \nicefrac{\pi}{2}$, yaitu sudut 90 derajat.
Galatnya $3.8\%$ ketika bermula dari $\nicefrac{\pi}{4}$, yaitu sudut
45 derajat. Pada sudut awal 5 derajat, galatnya hanya $0.048 \%$.

\begin{myfig}
\capstart
%original files nlin-T-vs-T0-abs nlin-T-vs-T0-relerr
\diffyincludegraphics{width=6.24in}{width=9in}{nlin-T-vs-T0-abs-relerr}{%
Dua plot. Pada plot kiri, T dan T sub linear diplot sebagai fungsi theta
sub 0 untuk theta sub 0 antara 0 dan pi per 2.
Plot T sub linear adalah garis horizontal pada 2 pi, sedangkan plot T
bermula pada 2 pi di kiri, tetapi perlahan naik hingga jauh di atas 7 di
kanan. Plot kanan menunjukkan galat relatif kedua besaran, yang bermula dari
nol ketika theta sub 0 sama dengan 0, lalu naik hingga sekitar 0.15 di kanan.}
\caption{Plot $T$ dan $T_{\text{linear}}$ dengan $\frac{g}{L} =
1$ (kiri), serta plot galat relatif
$\frac{T-T_{\text{linear}}}{T}$ (kanan), untuk $\theta_0$ antara 0 dan
$\nicefrac{\pi}{2}$. \label{fig:TvsT0}}
\end{myfig}

Walaupun tidak langsung terlihat dari rumus, benar bahwa
\begin{equation*}
\lim_{\theta_0 \uparrow \pi} T = \infty .
\end{equation*}
Artinya, periode menuju tak hingga ketika sudut awal mendekati titik
kesetimbangan tak stabil. Jadi, jika kita menempatkan bandul hampir terbalik,
bandul mungkin memerlukan waktu sangat lama untuk turun. Hal ini konsisten
dengan perilaku limitnya: bandul yang tepat terbalik tidak pernah melakukan
osilasi, sehingga keadaan itu dapat kita anggap memiliki periode tak hingga.

\subsection{Sistem pemangsa--mangsa atau Lotka--Volterra}

Salah satu penerapan sederhana sistem nonlinear yang paling umum adalah sistem
yang disebut \emph{\myindex{pemangsa--mangsa}} atau
\emph{\myindex{Lotka--Volterra}}%
\footnote{Dinamai menurut matematikawan, kimiawan, dan ahli statistika Amerika
\href{https://en.wikipedia.org/wiki/Alfred_J._Lotka}{Alfred James Lotka}
(1880--1949) dan matematikawan serta fisikawan Italia
\href{https://en.wikipedia.org/wiki/Vito_Volterra}{Vito Volterra}
(1860--1940).}
Sebagai contoh, sistem semacam ini muncul ketika dua spesies berinteraksi,
yang satu sebagai mangsa dan yang lain sebagai pemangsa. Karena itu, tidak
mengherankan bahwa persamaan-persamaan tersebut juga diterapkan dalam
ekonomi. Sistem ini juga muncul dalam reaksi kimia.
Dalam biologi, sistem persamaan ini menjelaskan variasi periodik alami pada
populasi berbagai spesies. Sebelum persamaan diferensial diterapkan, variasi
periodik pada populasi ini membingungkan para ahli biologi.

Kita tetap menggunakan contoh klasik kelinci dan rubah di sebuah hutan.
Contoh ini paling mudah dipahami.
\begin{equation*}
\begin{aligned}
& x = \# \text{ kelinci (mangsa),} \\
& y = \# \text{ rubah (pemangsa).}
\end{aligned}
\end{equation*}
Ketika kelinci berjumlah banyak, makanan bagi rubah melimpah sehingga
populasi rubah bertambah. Namun, ketika populasi rubah bertambah, lebih banyak
kelinci dimakan. Jadi, ketika rubah berjumlah banyak, populasi kelinci
seharusnya berkurang, dan sebaliknya.
Model Lotka--Volterra mengusulkan bahwa perilaku ini dideskripsikan oleh
sistem persamaan
\begin{equation*}
\begin{aligned}
& x' = (a-by)x, \\
& y' = (cx-d)y,
\end{aligned}
\end{equation*}
dengan $a,b,c,d$ sebagai sejumlah parameter yang mendeskripsikan interaksi
rubah dan kelinci\footnote{Interaksi ini tidak berakhir baik bagi kelinci.}.
Dalam model ini, semuanya merupakan bilangan positif.

Mari kita analisis gagasan di balik model ini. Model tersebut merupakan
gagasan yang sedikit lebih rumit dan didasarkan pada model populasi
eksponensial. Pertama, uraikan
\begin{equation*}
x' = (a-by)x = ax - byx .
\end{equation*}
Tanpa adanya rubah, kelinci diperkirakan tumbuh secara eksponensial.
Di sinilah suku $ax$ berperan: pertumbuhan populasi sebanding dengan populasi
itu sendiri. Kita mengasumsikan bahwa kelinci selalu menemukan cukup makanan
dan memiliki cukup ruang untuk berkembang biak. Namun, ada komponen lain
$-byx$, yakni populasi juga berkurang sebanding dengan jumlah rubah.
Keduanya dapat kita tulis sebagai $(a-by)x$, sehingga bentuknya menyerupai
pertumbuhan atau peluruhan eksponensial, tetapi konstantanya bergantung pada
jumlah rubah.

Persamaan bagi rubah sangat serupa. Uraikan lagi:
\begin{equation*}
y' = (cx-d)y = cxy-dy .
\end{equation*}
Rubah memerlukan makanan (kelinci) untuk berkembang biak: semakin banyak
makanan, semakin besar laju pertumbuhannya, sehingga muncul suku $cxy$.
Di sisi lain, terdapat kematian alami dalam populasi rubah, sehingga muncul
suku $-dy$.

Tanpa menunda lagi, mari kita mulai dengan sebuah contoh eksplisit. Misalkan
persamaannya adalah
\begin{equation*}
x' = (0.4-0.01y)x, \qquad y' = (0.003x-0.3)y .
\end{equation*}
Lihat \figurevref{fig:nlin-pred-prey} untuk potret fasenya. Dalam contoh ini,
masuk akal untuk memplot $x$ dan $y$ sebagai grafik terhadap waktu.
Karena itu, grafik kedua pada \figureref{fig:nlin-pred-prey} merupakan grafik
$x$ dan $y$ pada sumbu vertikal (mangsa $x$ adalah garis lebih tipis dengan
puncak lebih tinggi), terhadap waktu pada sumbu horizontal. Solusi khusus
yang digambar memiliki kondisi awal 20 rubah dan 50 kelinci.
\begin{myfig}
\capstart
%original files nlin-pred-prey-phase nlin-pred-prey-graphs
\diffyincludegraphics{width=6.24in}{width=9in}{nlin-pred-prey-phase-graphs}{%
Di sebelah kiri terdapat medan arah untuk x antara 0 dan 200 serta y antara
0 dan 100. Anak-anak panah berputar berlawanan arah jarum jam mengelilingi
titik kritis di (100,40). Beberapa lintasan diplot; bentuknya tidak benar-benar
melingkar. Lintasan bergerak kira-kira ke kanan sepanjang bagian bawah,
lalu berbelok naik kira-kira secara diagonal, kemudian turun kembali di sisi
kiri. Lintasan besar menyerupai segitiga yang sudut-sudutnya sangat
membulat. Di sebelah kanan terdapat plot untuk t antara 0 dan 40, dengan
grafik pemangsa berupa garis merah tebal dan mangsa berupa garis biru tipis.
Pada awalnya, jumlah pemangsa lebih rendah daripada jumlah mangsa.
Mula-mula mangsa bertambah dengan cepat dan mencapai puncak ketika pemangsa
sedang bertambah, lalu jumlah mangsa turun tajam. Ketika jumlah mangsa turun
hingga kira-kira setengah tinggi puncaknya, jumlah pemangsa mencapai puncak
dan mulai turun. Jumlah mangsa berhenti turun, pulih, lalu mulai bertambah
lagi, sementara jumlah pemangsa terus turun sedikit lebih lama. Kemudian
siklus ini berulang. Kurva mangsa jauh lebih tinggi, dengan puncak di atas
225, sedangkan pemangsa mencapai puncak sekitar 80. Nilai minimum pemangsa
sekitar 15, sedangkan nilai minimum mangsa sekitar 32. Kurva pemangsa lebih
curam ketika naik dan kurang curam ketika turun, sedangkan kurva mangsa naik
dan turun dengan kemiringan yang hampir sama, mungkin sedikit lebih lambat
ketika naik.}
\caption{Potret fase (kiri) dan grafik $x$ serta $y$ untuk
sebuah solusi contoh (kanan). \label{fig:nlin-pred-prey}}
\end{myfig}

Mari kita analisis apa yang tampak pada grafik. Kita bekerja dalam keadaan
umum alih-alih memasukkan bilangan tertentu. Kita mulai dengan mencari titik
kritis. Tetapkan $(a-by)x = 0$ dan $(cx-d)y = 0$.
Persamaan pertama dipenuhi jika $x=0$ atau $y=\nicefrac{a}{b}$. Jika $x=0$,
persamaan kedua mengakibatkan $y=0$. Jika $y= \nicefrac{a}{b}$, persamaan
kedua mengakibatkan $x=\nicefrac{d}{c}$.
Terdapat dua kesetimbangan: di $(0,0)$ ketika tidak ada hewan sama sekali,
dan di $(\nicefrac{d}{c},\nicefrac{a}{b})$.
Dalam contoh khusus kita, $x = \nicefrac{d}{c} = 100$ dan
$y = \nicefrac{a}{b} = 40$. Ini adalah titik dengan 100 kelinci dan 40 rubah.

Kita menghitung matriks Jacobian:
\begin{equation*}
\begin{bmatrix}
a-by & -bx \\
cy & cx-d
\end{bmatrix} .
\end{equation*}
Di titik asal $(0,0)$ kita memperoleh matriks
$\left[ \begin{smallmatrix}
a & 0 \\
0 & -d
\end{smallmatrix} \right]$, sehingga nilai eigennya adalah $a$ dan $-d$,
yang real dan berlainan tanda. Jadi, titik kritis di titik asal merupakan
pelana. Hal ini masuk akal. Jika kita bermula dengan sejumlah rubah tetapi
tanpa kelinci, rubah akan punah; artinya, kita mendekati titik asal.
Jika kita bermula tanpa rubah dan dengan beberapa kelinci, kelinci akan terus
berkembang biak tanpa kendali sehingga kita menjauhi titik asal.

Baiklah\@, bagaimana dengan titik kritis lainnya di
$(\nicefrac{d}{c},\nicefrac{a}{b})$? Di sini matriks Jacobian menjadi
\begin{equation*}
\begin{bmatrix}
0 & -\frac{bd}{c} \\
\frac{ac}{b} & 0
\end{bmatrix} .
\end{equation*}
Nilai eigen memenuhi $\lambda^2 + ad = 0$. Dengan kata lain,
$\lambda = \pm i \sqrt{ad}$. Karena nilai eigennya murni imajiner, kita
berada dalam kasus yang tidak dapat diputuskan hanya melalui linearisasi.
Kita mungkin memiliki pusat stabil, serapan spiral, atau sumber spiral.
Artinya, kesetimbangannya dapat stabil asimtotik, stabil, atau tak stabil.
Tentu saja, gambar di atas tampaknya menyiratkan bahwa titik itu merupakan
pusat stabil. Namun, jangan pernah hanya memercayai gambar. Mungkin
osilasinya semakin besar, tetapi hanya \emph{sangat} lambat. Tentu saja hal
ini buruk karena menyiratkan bahwa cepat atau lambat akan terjadi sesuatu
pada populasi kita. Selain itu, saya hanya menggambar contoh yang sangat
khusus dengan lintasan yang sangat khusus.

Bagaimana kita dapat memastikan bahwa kita berada dalam keadaan stabil?
Seperti telah kita katakan, untuk nilai eigen yang murni imajiner, kita perlu
bekerja sedikit lebih jauh. Sebelumnya kita menemukan bahwa, untuk sistem
konservatif, terdapat suatu besaran yang tetap konstan pada lintasan
sehingga lintasan harus membentuk gelung tertutup.
Kita dapat menggunakan teknik serupa di sini. Kita hanya perlu mencari
besaran apa yang tetap konstan. Setelah mencoba-coba, kita menemukan bahwa
konstanta
\begin{equation*}
C = \frac{y^a x^d}{e^{cx+by}} = y^a x^d e^{-cx-by}
\end{equation*}
tetap konstan. Besaran semacam ini disebut \emph{\myindex{konstanta
gerak}}. Mari kita periksa bahwa $C$ benar-benar merupakan konstanta gerak.
Bagaimana memeriksanya? Konstanta adalah sesuatu yang tidak berubah terhadap
waktu, jadi mari kita hitung turunannya terhadap waktu:
\begin{equation*}
C' = 
a y^{a-1}y' x^d e^{-cx-by}
+
y^a d x^{d-1} x' e^{-cx-by}
+
y^a x^d e^{-cx-by} (-cx'-by') .
\end{equation*}
Persamaan kita menyatakan nilai $x'$ dan $y'$, jadi mari kita substitusikan:
\begin{equation*}
\begin{split}
C' & = 
a y^{a-1} (cx-d)y x^d e^{-cx-by}
+
y^a d x^{d-1} (a-by)x e^{-cx-by}
\\
& \phantom{mm} +
y^a x^d e^{-cx-by} \bigl(-c(a-by)x-b(cx-d)y\bigr)
\\
& =
y^a x^d e^{-cx-by}
\Bigl(
a (cx-d)
+
d (a-by)
+
\bigl(-c(a-by)x-b(cx-d)y\bigr) \Bigr)
\\
& = 
0 .
\end{split}
\end{equation*}
Jadi, $C$ konstan sepanjang lintasan. Bahkan, bentuk
$C = \frac{y^a x^d}{e^{cx+by}}$ memberi kita persamaan implisit bagi
lintasan. Bagaimanapun juga, setelah menemukan konstanta gerak ini,
lintasan harus berupa kurva-kurva sederhana---yakni kurva level dari
$\frac{y^a x^d}{e^{cx+by}}$. Ternyata, titik kritis
$(\nicefrac{d}{c},\nicefrac{a}{b})$ merupakan titik maksimum bagi $C$
(dibiarkan sebagai latihan). Jadi, $(\nicefrac{d}{c},\nicefrac{a}{b})$
adalah titik kesetimbangan stabil, dan kita tidak perlu khawatir rubah serta
kelinci punah atau populasinya meledak.

Satu kekurangan model yang bagus ini ialah bahwa jumlah rubah dan kelinci
merupakan besaran diskret, sedangkan kita memodelkannya dengan variabel
kontinu. Model kita tidak mempersoalkan keberadaan, misalnya, 0.1 ekor rubah
di hutan, padahal kenyataannya hal itu tidak masuk akal. Hampirannya wajar
selama jumlah rubah dan kelinci besar, tetapi kurang masuk akal untuk jumlah
kecil. Kita harus berhati-hati ketika menafsirkan hasil apa pun dari model
semacam ini.

Konsekuensi menarik (dan mungkin berlawanan dengan intuisi) dari model ini
adalah bahwa menambahkan hewan ke hutan dapat menyebabkan kepunahan, sebab
variasinya akan menjadi terlalu besar dan salah satu populasi akan mendekati
nol. Sebagai contoh, misalkan terdapat 20 rubah dan 50 kelinci seperti
sebelumnya, tetapi sekarang kita memasukkan lebih banyak rubah sehingga
jumlahnya menjadi 200. Jika perhitungan dijalankan, kita mendapati bahwa
jumlah kelinci akan merosot menjadi sedikit lebih dari 1 ekor di seluruh
hutan. Dalam kenyataan, kemungkinan besar ini berarti kelinci punah, lalu
rubah juga akan punah karena tidak memiliki makanan.

Menunjukkan bahwa suatu sistem persamaan memiliki solusi stabil dapat menjadi
persoalan yang sangat sulit. Ketika Isaac Newton mengemukakan hukum-hukum
gerak planetnya, ia membuktikan bahwa satu planet yang mengorbit satu matahari
merupakan sistem stabil. Namun, tata surya dengan lebih dari 1 planet ternyata
sungguh sangat sulit. Bahkan, sistem semacam itu berperilaku kaotik (lihat
\sectionref{sec:chaos}), yang berarti bahwa perubahan kecil pada kondisi awal
menimbulkan hasil jangka panjang yang sangat berbeda. Dari eksperimen numerik dan
pengukuran, kita mengetahui bahwa Bumi tidak akan terlempar ke ruang
hampa atau menabrak Matahari, setidaknya selama beberapa juta tahun.
Namun, kita tidak mengetahui apa yang terjadi setelah itu.

\subsection{Latihan}

\begin{exercise}
Ambil \emph{\myindex{persamaan bandul nonlinear teredam}}
$\theta '' + \mu \theta' + (\nicefrac{g}{L})
\sin \theta = 0$ untuk suatu $\mu > 0$ (artinya, terdapat gesekan).
\begin{tasks}
\task
Untuk menyederhanakan, misalkan $\mu = 1$ dan $\nicefrac{g}{L} = 1$.
Tentukan dan klasifikasikan titik-titik kritisnya.
\task
Lakukan hal yang sama untuk sebarang $\mu > 0$ dan sebarang $g$
serta $L$, tetapi dengan redaman kecil, khususnya
$\mu^2 < 4(\nicefrac{g}{L})$.
\task
Jelaskan makna temuan Anda dan apakah hasilnya sesuai dengan apa yang
Anda perkirakan dalam kenyataan.
\end{tasks}
\end{exercise}

\begin{exercise}
Misalkan kelinci tidak tumbuh secara eksponensial, tetapi secara logistik.
Secara khusus, tinjau
\begin{equation*}
x' = (0.4-0.01y)x - \gamma x^2, \qquad y' = (0.003x-0.3)y .
\end{equation*}
Untuk dua nilai $\gamma$ berikut, tentukan dan klasifikasikan semua titik
kritis di kuadran positif, yakni untuk $x \geq 0$ dan $y \geq 0$.
Kemudian sketsakan diagram fasenya. Bahas implikasinya bagi perilaku populasi
dalam jangka panjang.
\begin{tasks}(2)
\task
$\gamma=0.001$, 
\task
$\gamma=0.01$.
\end{tasks}
\end{exercise}

\begin{exercise}
\leavevmode
\begin{tasks}
\task Misalkan $x$ dan $y$ adalah variabel positif. Tunjukkan bahwa
$\frac{y x}{e^{x+y}}$ mencapai maksimum di $(1,1)$.
\task Misalkan $a,b,c,d$ adalah konstanta positif, dan misalkan pula $x$ serta
$y$ adalah variabel positif. Tunjukkan bahwa
$\frac{y^a x^d}{e^{cx+by}}$ mencapai maksimum di
$(\nicefrac{d}{c},\nicefrac{a}{b})$.
\end{tasks}
\end{exercise}

\begin{exercise}
Untuk persamaan bandul, misalkan kita mengambil lintasan yang menghasilkan
gerak berputar penuh, misalnya
$\omega = \sqrt{\frac{2g}{L} \cos \theta
+ \frac{2g}{L} + \omega_0^2}$. Ini adalah lintasan dengan kecepatan sudut
terendah $\omega_0$. Tentukan bentuk integral bagi waktu yang diperlukan
bandul untuk melakukan satu putaran penuh.
\end{exercise}

\begin{exercise}[menantang]
Ambil bandul dan misalkan posisi awalnya adalah $\theta = 0$.
\begin{tasks}
\task
Tentukan rumus $\omega$ yang memberikan lintasan dengan syarat
awal $(0,\omega_0)$. Petunjuk: Cari tahu bagaimana $C$ harus dinyatakan dalam
$\omega_0$.
\task
Tentukan kecepatan sudut kritis $\omega_1$ sedemikian sehingga, untuk setiap
kecepatan sudut awal yang lebih tinggi, bandul akan terus berputar
mengelilingi porosnya, sedangkan untuk setiap kecepatan sudut awal yang lebih
rendah, bandul hanya akan berayun ke depan dan ke belakang.
Petunjuk: Ketika bandul tidak melewati titik teratas, rumus untuk
$\omega$ tidak terdefinisi bagi beberapa nilai $\theta$.
\task
Menurut Anda, apa yang terjadi jika kondisi awalnya adalah $(0,\omega_1)$,
yakni sudut awalnya 0 dan kecepatan sudut awalnya tepat $\omega_1$?
\end{tasks}
\end{exercise}

\setcounter{exercise}{100}

\begin{exercise}
Ambil persamaan bandul nonlinear teredam $\theta '' + \mu \theta' +
(\nicefrac{g}{L})
\sin \theta = 0$ untuk suatu $\mu > 0$ (artinya, terdapat gesekan).
Misalkan gesekannya besar, khususnya $\mu^2 > 4 (\nicefrac{g}{L})$.
\begin{tasks}
\task
Tentukan dan klasifikasikan titik-titik kritisnya.
\task
Jelaskan makna temuan Anda dan apakah hasilnya sesuai dengan apa yang
Anda perkirakan dalam kenyataan.
\end{tasks}
\end{exercise}
\exsol{%
a) Titik-titik kritisnya adalah $\omega=0$, $\theta=k\pi$ untuk setiap
bilangan bulat $k$. Ketika $k$ ganjil, kita memperoleh titik pelana.
Ketika $k$ genap, kita memperoleh serapan. \quad
b)~Temuan ini berarti bandul akan langsung menuju salah satu serapan,
misalnya $(0,0)$, dan tidak akan berayun ke depan dan ke belakang.
Gesekannya terlalu besar untuk memungkinkan osilasi, persis seperti sistem
massa--pegas yang teredam berlebih.
}

\begin{exercise}
Misalkan kita memiliki sistem pemangsa--mangsa dengan rubah yang juga dibunuh
pada laju konstan $h$ ($h$ rubah dibunuh per satuan waktu):
$x' = (a-by)x,$  $y' = (cx-d)y - h$.
\begin{tasks}
\task Tentukan titik-titik kritis dan matriks Jacobian sistem.
\task Masukkan konstanta $a=0.4$, $b=0.01$, $c=0.003$, $d=0.3$, $h=10$.
Analisis titik-titik kritisnya. Menurut Anda, apa maknanya bagi hutan itu?
\end{tasks}
\end{exercise}
\exsol{%
a) Dengan menyelesaikan persamaan untuk titik kritis, kita memperoleh
$(0,-\nicefrac{h}{d})$ dan
$(\frac{bh+ad}{ac},\frac{a}{b})$.
Matriks Jacobian di $(0,-\nicefrac{h}{d})$ adalah
$\left[
\begin{smallmatrix}
a+bh/d & 0 \\
-ch/d & -d
\end{smallmatrix}
\right]$
dengan nilai eigen $a+bh/d$ dan $-d$. Nilai eigennya real dengan tanda
berlawanan sehingga kita memperoleh pelana. (Namun, dalam penerapannya kita
hanya meninjau kuadran positif, sehingga titik kritis ini tidak relevan.)
Di $(\frac{bh+ad}{ac},\frac{a}{b})$ kita memperoleh matriks Jacobian
$\left[
\begin{smallmatrix}
0 & -\frac{b \left( b h+a d\right) }{a c}\\
\frac{a c}{b} & \frac{b h+a d}{a}-d
\end{smallmatrix}
\right]$.
b)~Untuk bilangan khusus yang diberikan, titik kritis kedua adalah
$(\frac{550}{3},40)$
dan matriksnya adalah
$\left[
\begin{smallmatrix}
0 & -11/6 \\
3/25 & 1/4
\end{smallmatrix}
\right]$, yang memiliki nilai eigen
$\frac{5\pm i \sqrt{327}}{40}$. Oleh karena itu, terdapat sumber spiral;
solusinya berpilin ke luar. Solusi itu pada akhirnya mengenai salah satu
sumbu, $x=0$ atau $y=0$, sehingga salah satu populasi di hutan akan punah.
}

\begin{exercise}[menantang]
Misalkan rubah tidak pernah mati. Artinya, kita memiliki sistem
$x' = (a-by)x,$ $y' = cxy$.
Tentukan titik-titik kritisnya dan perhatikan bahwa titik-titik itu tidak
terisolasi. Apa yang akan terjadi pada populasi di hutan jika populasi awal kedua spesies
bernilai positif? Petunjuk: Pikirkan konstanta geraknya.
\end{exercise}
\exsol{%
Titik-titik kritis berada pada garis $x=0$. Di kuadran positif, $y'$ selalu
positif sehingga populasi rubah selalu bertambah.
Konstanta geraknya adalah $C = y^ae^{-cx-by}$.
Untuk setiap $C$, kurva ini harus mengenai sumbu $y$ (mengapa?),
sehingga lintasan akan mendekati suatu titik pada sumbu $y$ dan jumlah
kelinci akan menuju nol.
}


%We have a conservative equation and so (exercise) the
%trajectories are given by
%\begin{equation*}
%\omega = \pm \sqrt{ \frac{2g}{L} \cos \theta + C} ,
%\end{equation*}
%for various values of $C$.  Let us figure out what $C$ corresponds to an
%initial condition $(0,\omega_0)$.  A little bit of thought tells us that
%such a $C = \omega_0^2 - \frac{2g}{L}$.  Taking just the top part of the
%trajectory we get
%\begin{equation*}
%\omega = \sqrt{ \frac{2g}{L} \cos \theta - \frac{2g}{L} + \omega_0^2} .
%\end{equation*}
%What we are trying to do is figure out when this will have no
%\myquote{gaps,} that
%is when what is under the square root is always positive.  The minimum is
%clearly taken when $\theta$ is an odd multiple of $\pi$, in this case we
%will get precisely zero when
%\begin{equation*}
%0 = \frac{2g}{L} \cos \pi - \frac{2g}{L} + \omega_0^2 ,
%\end{equation*}
%or in other words, solving for $\omega_0$ (and assuming it is positive) we
%have
%\begin{equation*}
%\omega_0 = 2 \sqrt{\frac{g}{L}} .
%\end{equation*}
%In the case we graphed, that is, when $\frac{g}{L} = 1$, then this magic
%$\omega_0 = 2$.  Notice that the trajectory that seems to go through the 
%saddle points goes through the point $(0,2)$.

%%%%%%%%%%%%%%%%%%%%%%%%%%%%%%%%%%%%%%%%%%%%%%%%%%%%%%%%%%%%%%%%%%%%%%%%%%%%%%

\sectionnewpage

\section{Siklus limit}
\label{limitcycles:section}

\sectionnotes{kurang dari 1 pertemuan\EPref{, dibahas dalam \S6.1 dan \S6.4
dalam \cite{EP}}\BDref{, \S9.7 dalam \cite{BD}}}

Untuk sistem nonlinear, lintasan tidak harus sekadar mendekati atau
meninggalkan satu titik. Lintasan bahkan dapat mendekati himpunan yang lebih
besar, misalnya lingkaran atau kurva tertutup lainnya.

\begin{example}
\emph{\myindex{Osilator Van der Pol}}\footnote{Dinamai menurut fisikawan
Belanda
\href{https://en.wikipedia.org/wiki/Balthasar_van_der_Pol}{Balthasar van der
Pol} (1889--1959).}
adalah persamaan berikut:
\begin{equation*}
x''-\mu(1-x^2) x' + x = 0,
\end{equation*}
dengan $\mu$ sebagai suatu konstanta positif. Osilator Van der Pol berasal
dari kajian rangkaian listrik, tetapi diterapkan dalam berbagai bidang seperti
biologi, seismologi, dan ilmu fisika lainnya.

Untuk menyederhanakan, kita menggunakan $\mu = 1$. Diagram fase diberikan
pada plot kiri dalam \figurevref{fig:nlin-van-der-fig}, dengan $y=x'$ seperti
biasa. Perhatikan bahwa lintasan tampak sangat cepat menetap pada suatu
kurva tertutup. Di sebelah kanan terdapat plot satu solusi untuk $t=0$ sampai
$t=30$ dengan kondisi awal $x(0) = 0.1$ dan $x'(0) = 0.1$. Solusi tersebut
dengan cepat menuju solusi periodik.
\begin{myfig}
\capstart
%original files nlin-van-der-phase nlin-van-der-plot1
\diffyincludegraphics{width=6.24in}{width=9in}{nlin-van-der-phase-plot1}{%
Di sebelah kiri terdapat bidang fase (plot arah) osilator Van der Pol, dengan
sumbu horizontal berupa x dan sumbu vertikal berupa turunan x prima.
Sumbu horizontal maupun vertikal membentang dari minus 4 sampai 4.
Pada plot ini, anak-anak panah tampak bergerak mengelilingi titik asal searah
jarum jam. Namun, ada satu lintasan yang merupakan kurva tertutup (sebuah
gelung), yang disebut siklus limit. Bentuknya secara sangat kasar menyerupai
persegi panjang membulat yang terbaring. Kurva itu kira-kira membentang antara
minus 2 dan 2 pada sumbu horizontal x serta antara minus 3 dan 3 pada sumbu
vertikal. Sebuah lintasan yang bermula di dalam siklus limit ini berpilin
ke luar menuju siklus limit dan semakin mendekatinya. Demikian pula, beberapa
lintasan yang bermula di luar mendekati siklus limit dari luar dan semakin
mendekatinya. Baik dari dalam maupun luar, pendekatannya sangat cepat, dan
lintasan bahkan belum satu kali mengelilinginya ketika sudah tidak dapat
dibedakan dari siklus limit. Di sebelah kanan terdapat plot pada bidang tx
untuk t antara 0 dan 30 dari sebuah solusi contoh yang bermula dekat titik
asal. Solusi itu dengan cepat memasuki osilasi periodik besar antara kira-kira
minus 2 dan 2.}
\caption{Potret fase (kiri) dan grafik sebuah solusi contoh
osilator Van der Pol.\label{fig:nlin-van-der-fig}}
\end{myfig}

Osilator Van der Pol merupakan contoh \emph{\myindex{osilasi relaksasi}}.
Kata relaksasi merujuk pada lonjakan tiba-tiba (bagian solusi yang sangat
curam). Untuk $\mu$ yang lebih besar, bagian curam ini lebih jelas.
Untuk $\mu$ yang lebih kecil, siklus limit lebih menyerupai lingkaran.
Bahkan, jika $\mu = 0$, maka $x''+x=0$, yang merupakan sistem linear dengan
sebuah pusat dan semua lintasannya berupa lingkaran.
\end{example}

Lintasan pada potret fase yang berupa kurva tertutup (kurva yang membentuk
gelung) disebut \emph{\myindex{lintasan tertutup}}.
\emph{\myindex{siklus limit}} adalah lintasan tertutup sedemikian sehingga
setidaknya satu lintasan lain berpilin masuk menuju lintasan tersebut
(atau berpilin keluar darinya). Sebagai contoh, kurva tertutup pada potret
fase persamaan Van der Pol adalah siklus limit.
Jika semua lintasan yang bermula dekat siklus limit berpilin masuk menuju
siklus tersebut, siklus limit disebut
\emph{stabil asimtotik}\index{siklus limit stabil asimtotik}.
Siklus limit pada osilator Van der Pol stabil asimtotik.

Untuk sebuah lintasan tertutup dalam sistem otonom, setiap solusi yang
bermula pada lintasan itu bersifat periodik. Kurva semacam ini disebut
\emph{\myindex{orbit periodik}}.
Lebih tepatnya, jika $\bigl(x(t),y(t)\bigr)$ merupakan solusi sedemikian
sehingga untuk suatu $t_0$ titik $\bigl(x(t_0),y(t_0)\bigr)$ terletak pada
orbit periodik, maka $x(t)$ maupun $y(t)$ merupakan fungsi periodik (dengan
periode yang sama). Artinya, terdapat suatu bilangan $P$ sedemikian sehingga
$x(t) = x(t+P)$ dan $y(t) = y(t+P)$.

Tinjau sistem
\begin{equation} \label{nlin:gensys}
x' = f(x,y), \qquad y' = g(x,y) ,
\end{equation}
dengan fungsi $f$ dan $g$ memiliki turunan kontinu pada suatu daerah $R$
di bidang.

\begin{theorem}[Poincar\'e--Bendixson\footnote{%
\href{https://en.wikipedia.org/wiki/Ivar_Otto_Bendixson}{Ivar Otto Bendixson}
(1861--1935) adalah seorang matematikawan Swedia.}]%
\index{Teorema Poincar\'e--Bendixson}
Misalkan $R$ adalah daerah tertutup dan terbatas (daerah pada bidang yang
mencakup batasnya dan jarak titik-titiknya dari titik asal terbatas), dan \eqref{nlin:gensys} tidak memiliki titik kritis di $R$.
Misalkan $\bigl(x(t), y(t)\bigr)$ merupakan solusi
\eqref{nlin:gensys} di $R$ yang ada untuk semua $t \geq t_0$. Maka solusi
tersebut merupakan fungsi periodik, atau solusi tersebut menuju suatu solusi
periodik di $R$.
\end{theorem}

Inti teorema ini adalah bahwa, jika Anda menemukan satu solusi yang ada untuk
semua $t$ yang cukup besar (yakni ketika $t$ menuju tak hingga) dan tetap
berada dalam daerah terbatas tanpa titik kritis, maka Anda telah menemukan
orbit periodik atau solusi yang berpilin menuju siklus limit---dalam jangka
panjang, solusi itu sangat dekat dengan fungsi periodik.
Jika $R$ memuat titik kritis, perilakunya mungkin lebih rumit: solusi dapat
menuju suatu gelung yang tersusun atas titik-titik kritis yang dihubungkan
oleh lintasan, atau solusi dapat sekadar menuju sebuah titik kritis.
Teorema tersebut lebih merupakan pernyataan kualitatif daripada alat yang
membantu perhitungan. Dalam praktiknya, sulit untuk menemukan solusi analitik
atau menunjukkan bahwa solusi tersebut ada untuk sepanjang waktu.
Namun, jika kita menduga solusinya ada, kita menyelesaikannya secara numerik
untuk waktu yang besar guna menghampiri siklus limit.
Keterbatasan lainnya ialah bahwa teorema ini hanya berlaku dalam dua dimensi.
Dalam tiga dimensi atau lebih, ruang geraknya terlalu luas.

Teorema ini berlaku bagi semua solusi osilator Van der Pol. Solusi yang
bermula di titik mana pun selain titik asal $(0,0)$ menuju solusi periodik
di sekitar siklus limit, sedangkan kondisi awal $(0,0)$ memberikan solusi
konstan $x=0$, $y=0$.

\begin{example}
Tinjau
\begin{equation*}
x' = y + {(x^2+y^2-1)}^2 x, \qquad
y' = -x + {(x^2+y^2-1)}^2 y.
\end{equation*}
Medan vektor dan solusi dengan kondisi awal
$(1.02,0)$, $(0.9,0)$, dan $(0.1,0)$ digambar pada
\figurevref{fig:nlin-unstable-limit-cycle}.

%\begin{mywrapfig}{3.25in}
\begin{myfig}
\capstart
\diffyincludegraphics{width=3in}{width=4.5in}{nlin-unstable-limit-cycle}{%
Plot arah suatu siklus limit tak stabil pada bidang xy, dengan siklus limit
berbentuk lingkaran satuan dan semua anak panah bergerak kira-kira searah
jarum jam. Sebuah lintasan dari dalam mendekati siklus limit, sedangkan
lintasan yang bermula tepat di luar siklus limit berpilin ke luar,
mengelilinginya satu setengah kali, lalu terlepas dan keluar dari gambar
melalui bagian atas.}
\caption{Contoh siklus limit tak stabil.\label{fig:nlin-unstable-limit-cycle}}
\end{myfig}
%\end{mywrapfig}

Perhatikan bahwa titik-titik pada lingkaran satuan (berjarak satu dari titik
asal) memenuhi $x^2+y^2-1=0$. Selain itu, $x(t) = \sin(t)$,
$y = \cos(t)$ merupakan solusi sistem. Oleh karena itu, kita memiliki
lintasan tertutup.
Untuk titik yang tidak terletak pada lingkaran satuan, suku kedua dalam $x'$ mendorong solusi
lebih jauh dari sumbu $y$ daripada pada sistem $x' = y$, $y' = -x$, sedangkan
$y'$ mendorong solusi lebih jauh dari sumbu $x$ daripada pada sistem linear
$x'=y$, $y' = -x$. Dengan kata lain, untuk semua kondisi awal lainnya,
lintasan akan berpilin ke luar.

Ini berarti bahwa, untuk kondisi awal di dalam lingkaran satuan, solusi
berpilin ke luar menuju solusi periodik pada lingkaran satuan, sedangkan untuk
kondisi awal di luar lingkaran satuan, solusi berpilin menjauh menuju tak
hingga. Oleh karena itu, lingkaran satuan merupakan siklus limit, tetapi tidak
stabil asimtotik.
Teorema Poincar\'e--Bendixson berlaku untuk titik awal di dalam lingkaran
satuan karena solusi-solusi tersebut tetap terbatas, tetapi tidak berlaku
untuk titik awal di luar sebab solusi-solusi tersebut menuju tak hingga.
\end{example}

Analisis serupa berlaku bagi sistem
\begin{equation*}
x' = y + {(x^2+y^2-1)} x, \qquad
y' = -x + {(x^2+y^2-1)} y.
\end{equation*}
Lingkaran satuan sekali lagi merupakan lintasan tertutup. Titik di luar
lingkaran satuan berpilin ke luar menuju tak hingga, tetapi sekarang titik
di dalam lingkaran satuan berpilin menuju titik kritis di titik asal.
Lingkaran satuan merupakan siklus limit yang tak stabil secara asimtotik,
dan semua lintasan berpilin menjauh darinya.
Jika kita meninjau sistem linear $x' = y$, $y' = -x$, semua lingkaran yang
berpusat di titik asal (bukan hanya lingkaran satuan) merupakan lintasan
tertutup, tetapi tidak ada lintasan yang berpilin menuju lintasan lain,
sehingga tidak satu pun lingkaran tersebut merupakan siklus limit.
%By Poincar\'e--Bendixson, 
Semua solusi ini periodik.

Menurut teorema Picard (\thmvref{sys:picardthm}), di setiap titik pada bidang
kita selalu dapat menemukan solusi \eqref{nlin:gensys} pada suatu selang waktu
singkat, baik ke depan maupun ke belakang, selama $f$ dan $g$ memiliki turunan kontinu.
Jadi, jika kita menemukan lintasan tertutup dalam sistem otonom, maka untuk
setiap titik awal di dalam lintasan tertutup itu, solusi akan ada untuk
sepanjang waktu dan tetap terbatas (solusi tetap berada di dalam lintasan
tertutup). Begitu kita menemukan solusi di atas yang mengelilingi lingkaran
satuan, kita mengetahui bahwa untuk setiap titik awal di dalam lingkaran,
solusi ada untuk sepanjang waktu dan teorema Poincar\'e--Bendixson berlaku.

\medskip

Selanjutnya, kita mencari kondisi ketika siklus limit (atau orbit periodik)
tidak ada. Kita mengasumsikan persamaan \eqref{nlin:gensys} terdefinisi pada
\emph{\myindex{daerah terhubung sederhana}}, yaitu daerah tanpa lubang yang
dapat kita kelilingi. Sebagai contoh, seluruh bidang merupakan daerah
terhubung sederhana. Demikian pula bagian dalam cakram satuan. Namun, seluruh
bidang dengan satu titik dihilangkan bukan merupakan daerah terhubung
sederhana karena memiliki \myquote{lubang} di titik asal.

\begin{theorem}[Bendixson--Dulac%
\footnote{%
\href{https://en.wikipedia.org/wiki/Henri_Dulac}{Henri Dulac} (1870--1955)
adalah seorang matematikawan Prancis.}]%
\index{Teorema Bendixson--Dulac}
Misalkan $R$ adalah daerah terhubung sederhana, dan ekspresi%
%Suppose $f$ and $g$ are functions with continuous derivatives in
%a simply connected region $R$, and
%If the expression%
\footnote{Biasanya ekspresi dalam Teorema Bendixson--Dulac adalah
$\frac{\partial (\varphi f)}{\partial x} + \frac{\partial (\varphi
g)}{\partial y}$
untuk suatu fungsi diferensiabel secara kontinu $\varphi$.
Untuk menyederhanakan,
kita hanya meninjau kasus $\varphi = 1$.}
\begin{equation*}
\frac{\partial f}{\partial x} + \frac{\partial g}{\partial y}
\end{equation*}
selalu positif atau selalu negatif pada $R$ (kecuali mungkin pada himpunan
kecil seperti titik-titik terisolasi atau kurva-kurva), maka sistem
\eqref{nlin:gensys} tidak memiliki lintasan tertutup di dalam $R$.
\end{theorem}

Teorema ini memberikan cara untuk menolak keberadaan lintasan tertutup,
dan karena itu juga cara untuk menolak keberadaan siklus limit.
Ekspresi $\frac{\partial f}{\partial x} + \frac{\partial g}{\partial y}$
mungkin tampak muncul begitu saja, tetapi ekspresi itu disebut
\emph{\myindex{divergensi}} medan vektor.
Divergensi mengukur seberapa besar medan vektor \myquote{mengembang} di setiap
titik dan muncul dalam banyak konteks lain.
Teorema tersebut menyatakan bahwa, jika medan vektor mengembang di seluruh
$R$ atau menyusut di seluruh $R$, maka tidak ada lintasan tertutup dan
karena itu tidak ada siklus limit.
Mungkin ini terasa intuitif---jika partikel bergerak mengikuti medan vektor
dan semakin menjauh satu sama lain, kita tidak mengharapkan partikel apa pun
bergerak dalam gelung.
Pengecualian mengenai titik atau kurva berarti bahwa divergensi dapat bernilai
nol di beberapa titik atau pada sebuah kurva, tetapi tidak pada himpunan yang
lebih besar.

\begin{example}
Tinjau $x'=y+y^2e^x$, $y'=x$ pada seluruh bidang (lihat
\examplevref{example:nlin-withexp}).
Bidang tersebut terhubung sederhana sehingga teorema dapat diterapkan.
Kita menghitung
$\frac{\partial f}{\partial x} + \frac{\partial g}{\partial y} =
y^2e^x+ 0$. Fungsi $y^2e^x$ selalu positif kecuali pada garis $y=0$.
Oleh karena itu, menurut teorema tersebut, sistem tidak memiliki lintasan
tertutup.
\end{example}

Salah satu cara informal yang umum, tetapi tidak sepenuhnya benar, untuk
menyatakan teorema tersebut ialah menyimpulkan bahwa tidak ada solusi periodik
yang tetap berada dalam $R$. Contoh di atas memiliki dua titik kritis dan
karena itu memiliki solusi konstan. Fungsi konstan bersifat periodik.
Kesimpulan teorema adalah bahwa tidak ada lintasan yang membentuk kurva
tertutup, atau dengan kata lain, tidak ada solusi periodik tak konstan yang
tetap berada dalam $R$.

\begin{example}
Tinjau sebuah contoh yang agak lebih rumit.
Ambil sistem $x'=-y-x^2$, $y'=-x+y^2$ (lihat
\examplevref{example:nlin-xplusy}). Kita menghitung
$\frac{\partial f}{\partial x} + \frac{\partial g}{\partial y} =
-2x + 2y=2(-x+y)$. Ekspresi ini dapat bernilai positif maupun negatif, sehingga jika kita
membicarakan seluruh bidang, teorema tidak dapat langsung diterapkan.
Namun, kita dapat menerapkannya pada himpunan tempat $-x+y \geq 0$.
Menurut teorema, tidak ada lintasan tertutup dalam himpunan tersebut.
Demikian pula, tidak ada lintasan tertutup dalam himpunan
$-x+y \leq 0$. Kita belum dapat menyimpulkan bahwa tidak ada lintasan
tertutup di seluruh bidang. Sejauh yang kita ketahui, mungkin separuh
lintasan itu berada dalam himpunan $-x+y \geq 0$ dan separuh lainnya dalam
himpunan $-x+y \leq 0$.

Kuncinya adalah meninjau garis tempat $-x+y=0$, atau $x=y$.
Pada garis ini,
$x' = -y-x^2 = -x-x^2$ dan $y' = -x+y^2 = -x+x^2$.
Secara khusus, jika $x=y$, maka $x' \leq y'$. Jadi anak-anak panah, yakni
vektor-vektor $(x',y')$, selalu menunjuk ke dalam himpunan tempat
$-x+y \geq 0$. Kita tidak mungkin bermula dalam himpunan
$-x+y \geq 0$ lalu masuk ke himpunan $-x+y \leq 0$. Begitu berada dalam
himpunan $-x+y \geq 0$, kita tetap di sana. Jadi, tidak ada lintasan
tertutup yang dapat memiliki titik dalam kedua himpunan.
\end{example}

\begin{example}
Tinjau
$x' = y+(x^2+y^2-1)x$,
$y' = -x +(x^2+y^2-1)y$ pada daerah $R$ yang diberikan oleh
$x^2+y^2 > \frac{1}{2}$. Artinya, $R$ adalah daerah di luar lingkaran
berjari-jari $\frac{1}{\sqrt{2}}$ yang berpusat di titik asal. Terdapat
lintasan tertutup dalam $R$, yaitu $x=\cos(t)$, $y=\sin(t)$.
Selain itu,
\begin{equation*}
\frac{\partial f}{\partial x} + 
\frac{\partial g}{\partial y} = 4x^2+4y^2-2 ,
\end{equation*}
yang selalu positif pada $R$. Jadi, apa yang terjadi? Teorema
Bendixson--Dulac tidak berlaku karena daerah $R$ tidak terhubung
sederhana---daerah tersebut memiliki lubang, yaitu lingkaran yang kita
hilangkan!
\end{example}

\subsection{Latihan}

\begin{exercise}
Tunjukkan bahwa sistem-sistem berikut tidak memiliki lintasan tertutup.
\begin{tasks}(2)
\task $x'=x^3+y,\quad y'=y^3+x^2$,
\task $x'=e^{x-y},\quad y'=e^{x+y}$,
\task $x'=x+3y^2-y^3,\quad y'=y^3+x^2$.
\end{tasks}
\end{exercise}

\begin{exercise}
Rumuskan syarat agar sistem linear berukuran 2 kali 2
${\vec{x}}' = A \vec{x}$ bukan merupakan pusat dengan menggunakan
teorema Bendixson--Dulac. Artinya, teorema tersebut menyatakan sesuatu
tentang unsur-unsur tertentu dari $A$.
\end{exercise}

\begin{exercise}
Jelaskan mengapa Teorema Bendixson--Dulac tidak berlaku bagi sistem
konservatif apa pun $x''+h(x) = 0$.
\end{exercise}

\begin{exercise}
Sistem seperti $x'=x, y'=y$ memiliki solusi yang ada untuk sepanjang waktu
$t$, tetapi tidak memiliki lintasan tertutup. Jelaskan mengapa Teorema
Poincar\'e--Bendixson tidak berlaku.
\end{exercise}

\begin{exercise}
Persamaan diferensial juga dapat diberikan dalam sistem koordinat yang
berbeda. Misalkan kita memiliki sistem $r' = 1-r^2$, $\theta' = 1$ yang
diberikan dalam koordinat polar. Tentukan semua lintasan tertutup dan
periksa apakah lintasan tersebut merupakan siklus limit; jika ya, tentukan
apakah lintasan itu stabil asimtotik.
\end{exercise}


\setcounter{exercise}{100}

\begin{exercise}
Tunjukkan bahwa sistem-sistem berikut tidak memiliki lintasan tertutup.
\begin{tasks}(2)
\task $x'=x+y^2,\quad y'=y+x^2$,
\task $x'=-x\sin^2(y),\quad y'=e^x$,
\task $x'=xy^2,\quad y'=x+x^2$.
\end{tasks}
\end{exercise}
\exsol{%
Gunakan Teorema Bendixson--Dulac.
a) $f_x+g_y = 1+1 > 0$, sehingga tidak ada lintasan tertutup.
b) $f_x+g_y = -\sin^2(y)+0 < 0$ untuk semua $x,y$, kecuali garis-garis yang
diberikan oleh $y=k\pi$ (tempat kita memperoleh nol), sehingga tidak ada
lintasan tertutup.
c) $f_x+g_y = y^2 + 0 > 0$ untuk semua $x,y$, kecuali garis yang diberikan
oleh $y=0$ (tempat kita memperoleh nol), sehingga tidak ada lintasan
tertutup.
}

\begin{exercise}
Misalkan suatu sistem otonom pada bidang memiliki solusi
$x=\cos(t)(1+e^{-t})$, $y=\sin(t)(1+e^{-t})$. Apa yang dapat Anda katakan
tentang sistem tersebut (khususnya tentang siklus limit dan solusi periodik)?
\end{exercise}
\exsol{%
Dengan menggunakan Teorema Poincar\'e--Bendixson, sistem memiliki sebuah
siklus limit, yaitu lingkaran satuan yang berpusat di titik asal, karena
$x=\cos(t)(1+e^{-t})$, $y=\sin(t)(1+e^{-t})$ semakin mendekati lingkaran
satuan. Jadi, $x=\cos(t)$, $y=\sin(t)$ merupakan solusi periodiknya.
}

\begin{exercise}
Tunjukkan bahwa siklus limit osilator Van der Pol (untuk $\mu > 0$)
tidak mungkin seluruhnya terletak dalam himpunan tempat $-1 < x < 1$.
Bandingkan dengan \figurevref{fig:nlin-van-der-fig}.
\end{exercise}
\exsol{%
$f(x,y) = y$, $g(x,y) = \mu(1-x^2)y-x$. Jadi,
$f_x+g_y = \mu(1-x^2)$. Teorema Bendixson--Dulac menyatakan bahwa tidak
ada lintasan tertutup yang seluruhnya terletak dalam himpunan $x^2 < 1$.
}

\begin{exercise}
Tinjau sistem $r' = \sin(r)$, $\theta' = 1$ yang diberikan dalam koordinat
polar. Tentukan semua lintasan tertutup.
\end{exercise}
\exsol{%
Lintasan tertutup adalah lintasan tempat $\sin(r) = 0$. Oleh karena itu,
semua lingkaran yang berpusat di titik asal dan berjari-jari kelipatan positif
$\pi$ merupakan lintasan tertutup.
}


%%%%%%%%%%%%%%%%%%%%%%%%%%%%%%%%%%%%%%%%%%%%%%%%%%%%%%%%%%%%%%%%%%%%%%%%%%%%%%

\sectionnewpage
\section{Chaos} \label{sec:chaos}

%mbxINTROSUBSECTION

\sectionnotes{1 pertemuan\EPref{, \S6.5 dalam \cite{EP}}\BDref{,
\S9.8 dalam \cite{BD}}}

Anda tentu pernah mendengar cerita tentang kepakan sayap kupu-kupu di
Amazon yang menyebabkan badai di Atlantik Utara.  Pada bagian sebelumnya,
kita menyebutkan bahwa perubahan kecil pada kondisi awal planet-planet dapat
menghasilkan konfigurasi planet yang sangat berbeda dalam jangka panjang.
Ini merupakan contoh
\emph{\myindex{sistem kaotik}}.
Chaos matematis sebenarnya bukanlah kekacauan---ada tatanan yang persis di
balik apa yang tampak.  Semuanya tetap deterministik.  Namun, sistem kaotik
sangat peka terhadap kondisi awal.  Ini juga berarti bahwa galat kecil yang
ditimbulkan oleh hampiran numerik akan dengan sangat cepat menjadi galat
besar, sehingga penghampiran numerik untuk jangka waktu panjang hampir tidak
mungkin dilakukan.  Hal ini menjadi bagian besar dari kesulitannya karena,
secara umum, sistem kaotik tidak dapat diselesaikan secara analitik.

Ambillah cuaca sebagai contoh, mungkin sistem kaotik yang paling dikenal.
Perubahan kecil pada kondisi awal
(misalnya suhu di setiap titik atmosfer) menghasilkan prediksi yang sangat
berbeda dalam waktu relatif singkat, sehingga kita tidak dapat meramalkan
cuaca dengan akurat.  Lagi pula, kita sebenarnya tidak mengetahui syarat
awalnya secara tepat.  Kita mengukur suhu di beberapa titik dengan sejumlah
galat, lalu memperkirakan entah bagaimana keadaan di antara titik-titik itu.
Tidak mungkin kita dapat mengukur secara akurat pengaruh setiap kepakan sayap
kupu-kupu.  Selanjutnya kita menyelesaikan persamaannya secara numerik dan
dengan demikian memperkenalkan galat baru.  Anda sebaiknya tidak mengandalkan
ramalan cuaca untuk waktu lebih dari beberapa hari ke depan.  Namun, kita akan
melihat bahwa masih ada informasi yang dapat diperoleh tentang sistem kaotik
pada skala waktu yang lebih panjang.  Dalam konteks cuaca, kita dapat
mempelajari iklim.

Perilaku kaotik pertama kali diperhatikan oleh Edward
Lorenz\footnote{\href{https://en.wikipedia.org/wiki/Edward_Norton_Lorenz}{Edward Norton Lorenz} (1917--2008) adalah
matematikawan dan meteorolog Amerika.}
pada tahun 1960-an ketika ia mencoba memodelkan konveksi (gerakan) udara yang
ditimbulkan oleh panas.  Lorenz menelaah sistem yang relatif sederhana:
\begin{equation*}
x' = -10x +10y, \qquad y' = 28x-y-xz, \qquad z'=-\frac{8}{3}z + xy .
\end{equation*}
Perubahan kecil pada kondisi awal menghasilkan solusi yang sangat berbeda
setelah waktu yang cukup singkat.

\begin{mywrapfigsimp}{0.95in}{1.25in}
\noindent
\inputpdft{chaos-pend}{%
Diagram sebuah bandul yang dipasang pada ujung bandul lain.  Massa pada
kedua ujung digambarkan dengan ukuran berbeda, dan kedua panjangnya juga
berbeda.}
\end{mywrapfigsimp}
Contoh sederhana yang dapat dicoba oleh pembaca dan menunjukkan perilaku
kaotik adalah bandul ganda.  Persamaan untuk susunan ini agak rumit dan
penurunannya cukup
\myindex{menjemukan}, sehingga kita tidak akan menuliskannya.  Gagasannya
adalah memasang sebuah bandul pada ujung bandul lain.  Gerakan massa di bagian
bawah akan tampak kaotik.  Jenis sistem kaotik ini menjadi dasar sejumlah
mainan meja kantor yang unik.  Versi sederhananya mudah dibuat.  Ambillah
seutas tali.  Ikatkan dua mur yang berat pada dua titik berbeda di tali;
satu di ujung, dan satu lagi agak di atasnya.  Sekarang beri mur bawah sedikit
dorongan.  Selama ayunannya tidak terlalu besar dan tali tetap tegang, Anda
memiliki sistem bandul ganda.

\subsection{Persamaan Duffing dan atraktor aneh}

Mari kita pelajari apa yang disebut \emph{\myindex{persamaan Duffing}}:
\begin{equation*}
x'' + a x' + bx + cx^3 = C \cos(\omega t) .
\end{equation*}
Di sini $a$, $b$, $c$, $C$, dan $\omega$ adalah konstanta.
Selain suku $c x^3$, persamaan ini tampak seperti sistem massa-pegas yang
dipaksa.  Adanya $c x^3$ berarti pegas tersebut tidak sepenuhnya mematuhi
hukum Hooke (dan memang tidak ada pegas nyata yang mematuhinya secara persis).
Ketika $c$ tidak nol, persamaan ini tidak mempunyai solusi bentuk tertutup,
sehingga kita harus menggunakan solusi numerik, sebagaimana lazimnya untuk
sistem nonlinear.  Tidak semua pilihan konstanta dan kondisi awal menunjukkan
perilaku kaotik.  Mari kita pelajari
\begin{equation*}
x''+0.05 x' + x^3 = 8\cos(t) .
\end{equation*}

Persamaan ini tidak otonom, sehingga kita tidak dapat menggambar medan vektor
pada bidang fase.  Kita tetap dapat menggambar lintasan, dengan $y=x'$ seperti
biasa.  Dalam \figurevref{nlin:duf-two-traj}, kita menggambar lintasan untuk
$t$ dari 0 hingga 15 bagi dua kondisi awal yang sangat berdekatan, yaitu
$(2,3)$ dan $(2,2.9)$, serta menggambar solusinya di ruang $(x,t)$.  Pada
awalnya kedua lintasan berdekatan, tetapi tidak lama kemudian keduanya
menyimpang secara nyata.  Kepekaan terhadap kondisi awal inilah yang kita
maksud ketika mengatakan bahwa sistem tersebut berperilaku kaotik.

\begin{myfig}
\capstart
%original files nlin-duf-two-traj nlin-duf-two-sols
\diffyincludegraphics{width=6.24in}{width=9in}{nlin-duf-two-traj-sols}{%
Di sebelah kiri terdapat grafik dua lintasan dalam ruang fase (ruang dengan
x pada sumbu mendatar dan x prima pada sumbu tegak).  Kedua lintasan bermula
di sebelah kanan grafik dengan posisi yang sangat berdekatan dan mengikuti
jalur yang mirip selama beberapa saat, tetapi kemudian menyimpang dengan
sangat cepat dan akhirnya tidak lagi dapat dikenali sebagai pasangan.  Bentuk
yang diikutinya kira-kira berupa gelung besar di sebelah kiri lalu gelung yang
lebih kecil di sebelah kanan, dan kadang-kadang membentuk gelung lebih kecil
di antaranya.  Di sebelah kanan terdapat grafik x terhadap t dari 0 hingga 15
untuk kedua solusi, yang naik turun ketika mengikuti bentuk-bentuk tersebut.
Kedua grafik hampir tidak dapat dibedakan di sisi kiri dan hampir sama hingga
kira-kira t sama dengan 7, tetapi ketika t mencapai 15 keduanya tampak berbeda
dan terpisah cukup jauh.}
\caption{Di kiri, dua lintasan dalam ruang fase untuk $0 \leq t \leq 15$ bagi persamaan Duffing,
yang satu dengan kondisi awal $(2,3)$ dan yang lain dengan kondisi awal $(2,2.9)$.  Di
kanan, kedua solusi dalam ruang $(x,t)$. \label{nlin:duf-two-traj}}
\end{myfig}

\begin{myfig}
\capstart
\diffyincludegraphics{width=6.24in}{width=9in}{nlin-duf-long}{%
Grafik lebar untuk t antara 0 dan 100, dengan sumbu x tegak yang membentang
kira-kira dari -3 hingga 3.  Grafik menunjukkan fungsi yang naik turun
berkali-kali dengan irama umum yang kurang lebih sama; puncak dan lembahnya
secara sangat kasar mempunyai tinggi yang sama, tetapi setiap puncak berbentuk
sangat berbeda dan memiliki lekukan khasnya sendiri.}
\caption{Solusi persamaan Duffing yang diberikan untuk $t$ dari 0 hingga 100.
\label{nlin:duf-long}}
\end{myfig}

\pagebreak[2]
Mari kita tinjau perilaku jangka panjangnya.
Dalam \figurevref{nlin:duf-long}, kita menggambar solusi $x$ dengan kondisi awal
$(2,3)$ untuk rentang waktu yang lebih panjang.  Sulit melihat pola apa pun
dalam bentuk solusi tersebut selain bahwa solusi itu tampak berosilasi, tetapi
setiap osilasi tampak cukup khas.  Osilasi memang diharapkan karena adanya
suku pemaksa.  Perlu disebutkan bahwa, agar gambar itu dihasilkan dengan
akurat, algoritme numeriknya menggunakan jumlah langkah yang luar biasa
besar\footnote{Sebagai rujukan, sebenarnya digunakan 30.000 langkah dengan
algoritme Runge--Kutta; lihat latihan dalam \sectionref{numer:section}.},
sebab dalam sistem kaotik bahkan galat kecil pun merambat dengan cepat.


Sistem kaotik sangat sulit dianalisis, begitu pula tatanan di balik apa yang
tampak kacau sangat sulit ditemukan.  Namun, mari kita coba melakukan sesuatu
yang dahulu kita lakukan untuk sistem massa-pegas standar.  Salah satu cara
kita menganalisis sistem itu adalah dengan menentukan perilaku jangka
panjangnya (yang tidak bergantung pada kondisi awal).  Dari gambar di atas jelas
bahwa kita tidak akan memperoleh uraian eksak yang rapi mengenai perilaku
jangka panjang sistem kaotik ini.  Akan tetapi, barangkali kita dapat menemukan
suatu keteraturan pada apa yang terjadi selama setiap \myquote{osilasi} dan
pada kesamaan di antara osilasi-osilasi tersebut.

Konsep yang akan kita bahas adalah \emph{\myindex{penampang Poincar\'e}}%
\footnote{Dinamai menurut cendekiawan serbabisa asal Prancis
\href{https://en.wikipedia.org/wiki/Henri_Poincar\%C3\%A9}{Jules Henri
Poincar\'e} (1854--1912).}.  Alih-alih melihat $t$ dalam suatu interval,
kita melihat kedudukan sistem pada suatu urutan waktu tertentu.  Bayangkan
menyalakan lampu strobo dengan frekuensi tetap dan menggambar titik tempat
solusi berada pada setiap kilatan.  Frekuensi strobo yang tepat bergantung
pada sistem yang sedang ditinjau.  Untuk persamaan Duffing paksa (dan
sistem serupa lainnya), frekuensi yang tepat adalah frekuensi suku pemaksa.
Untuk persamaan Duffing di atas, carilah solusi
$\bigl(x(t),y(t)\bigr)$, lalu perhatikan titik-titik
\begin{equation*}
\bigl(x(0),y(0)\bigr), \quad
\bigl(x(2\pi),y(2\pi)\bigr), \quad
\bigl(x(4\pi),y(4\pi)\bigr), \quad
\bigl(x(6\pi),y(6\pi)\bigr), \quad \ldots
\end{equation*}
Karena kita sebenarnya tidak berminat pada bagian transien solusi, yaitu
bagian solusi yang bergantung pada kondisi awal, kita melewatkan sejumlah
langkah di awal.  Misalnya, kita dapat melewatkan 100 langkah pertama dan
mulai menggambar titik pada $t = 100(2\pi)$, yaitu
\begin{equation*}
\bigl(x(200\pi),y(200\pi)\bigr), \quad
\bigl(x(202\pi),y(202\pi)\bigr), \quad
\bigl(x(204\pi),y(204\pi)\bigr), \quad \ldots
\end{equation*}
Gambar titik-titik ini disebut penampang Poincar\'e.  Setelah cukup banyak
titik digambar, muncullah pola yang menarik dalam
\figurevref{nlin:strange} (gambar sebelah kiri), yang disebut
\emph{\myindex{atraktor aneh}}.

\begin{myfig}
\capstart
%original files nlin-strange nlin-strange2
\diffyincludegraphics{width=6.24in}{width=9in}{nlin-strange-strange2}{%
Dua grafik ribuan titik pada atraktor aneh.  Keduanya terdeformasi dengan
cara yang sedikit berbeda karena grafik sebelah kanan mengalami pergeseran
fase.  Titik-titik itu tampak mengendap pada pita-pita tipis yang melengkung,
saling membelit, dan terdistorsi seolah-olah dilihat melalui cermin distorsi
di wahana rumah cermin.}
\caption{Atraktor aneh.  Grafik kiri tidak mengalami pergeseran fase,
sedangkan grafik kanan mengalami pergeseran fase
$\nicefrac{\pi}{4}$. \label{nlin:strange}}
\end{myfig}

Untuk suatu barisan titik, \emph{\myindex{atraktor}} adalah himpunan yang
semakin lama semakin didekati oleh titik-titik dalam barisan tersebut; dengan
kata lain, titik-titik itu tertarik kepadanya.  Penampang Poincar\'e sebenarnya
bukan atraktor itu sendiri, tetapi karena titik-titiknya sangat dekat dengan
atraktor, kita dapat melihat bentuknya.  Atraktor aneh merupakan himpunan yang
sangat rumit.  Himpunan ini memiliki struktur fraktal: sejauh apa pun Anda
memperbesar gambar, Anda akan terus melihat struktur rumit yang sama.

Kondisi awal tidak berpengaruh pada hasil akhirnya.  Jika kita memulai dengan
kondisi awal lain, titik-titik pada akhirnya bergerak mendekati atraktor.
Karena itu, selama beberapa titik pertama dibuang, kita memperoleh gambar yang
sama.  Demikian pula, galat kecil dalam hampiran numerik tidak menjadi masalah
di sini.

Hal yang menakjubkan ialah bahwa sistem kaotik seperti persamaan Duffing sama
sekali tidak acak.  Ada tatanan yang sangat rumit di dalamnya, dan atraktor
aneh mengungkapkan sesuatu tentang tatanan ini.  Kita tidak dapat menyatakan
dengan tepat keadaan akhir sistem, tetapi dengan frekuensi strobo yang tetap,
kita dapat mempersempit kemungkinannya menjadi titik-titik pada atraktor.

Jika kita menggunakan pergeseran fase, misalnya $\nicefrac{\pi}{4}$, dan
memperhatikan waktu-waktu
\begin{equation*}
\nicefrac{\pi}{4}, \quad
2\pi+\nicefrac{\pi}{4}, \quad
4\pi+\nicefrac{\pi}{4}, \quad
6\pi+\nicefrac{\pi}{4}, \quad
\ldots
\end{equation*}
kita memperoleh atraktor yang sedikit berbeda.  Gambarnya adalah bagian kanan
\figurevref{nlin:strange}.  Atraktor itu tampak seperti atraktor semula yang
diputar, digeser, dan sedikit didistorsi.  Untuk setiap pergeseran fase, Anda
dapat menemukan himpunan titik yang secara berkala terus didekati kembali oleh
sistem.

Pelajari gambar-gambarnya dan terutama perhatikan skalanya---di mana atraktor
tersebut berada pada bidang fase.  Perhatikan daerah tempat atraktor aneh itu
berada dan bandingkan dengan grafik lintasan dalam
\figurevref{nlin:duf-two-traj}.

Mari kita bandingkan bagian ini dengan pembahasan osilasi paksa dalam
\sectionref{forcedo:section}.  Tinjau
\begin{equation*}
x''+2p x' + \omega_0^2 x = \frac{F_0}{m} \cos (\omega t) .
\end{equation*}
Persamaan ini seperti persamaan Duffing, tetapi tanpa suku $x^3$.
Solusi periodik tunaknya berbentuk
\begin{equation*}
x = C \cos (\omega t + \gamma) .
\end{equation*}
Dengan pengamatan strobo berfrekuensi $\omega$, kita memperoleh satu titik
dalam ruang fase.  Atraktor dalam keadaan ini adalah satu titik---hasil yang
wajar karena sistemnya tidak kaotik.  Sistem ini justru merupakan kebalikan
dari sistem kaotik: setiap perbedaan yang ditimbulkan oleh kondisi awal lenyap
dengan sangat cepat, dan sistem selalu menetap pada gerak periodik tunak yang
sama.

\subsection{Sistem Lorenz}

Untuk menemukan perilaku kaotik dalam dua dimensi, kita harus mempelajari
sistem yang dipaksa atau sistem nonotonom, seperti persamaan Duffing.
Teorema Poincar\'e--Bendixson menyatakan bahwa solusi suatu sistem otonom dua
dimensi yang ada untuk seluruh waktu mendatang akan menuju tak hingga,
bersifat periodik, mendekati suatu gelung, atau mungkin mendekati titik
kritis.  Itu sama sekali bukan perilaku kaotik yang kita cari.

Dalam tiga dimensi, bahkan sistem otonom dapat bersifat kaotik.  Mari kita
kembali sejenak ke \myindex{sistem Lorenz}
\begin{equation*}
x' = -10x +10y, \qquad y' = 28x-y-xz, \qquad z'=-\frac{8}{3}z + xy .
\end{equation*}
Sistem Lorenz adalah sistem otonom tiga dimensi yang menunjukkan perilaku
kaotik.  Lihat \figurevref{nlin:lorenz} untuk contoh lintasan, yang kini berupa
kurva dalam ruang tiga dimensi.
\begin{myfig}
\capstart
\diffyincludegraphics{width=3in}{width=4.5in}{nlin-lorenz}{%
Sebuah lintasan dalam ruang xyz tiga dimensi.  Bentuknya kira-kira dua gelung
yang bertemu pada suatu sudut; lintasan mengitari salah satu gelung satu kali
atau lebih, lalu tampak berpindah secara acak ke gelung lain, dan kemudian
kembali lagi setelah satu putaran atau lebih.  Setiap putaran memiliki
jari-jari yang berbeda.  Sebagian orang menggambarkan bentuk ini seperti
kupu-kupu.}
\caption{Sebuah lintasan dalam sistem Lorenz. \label{nlin:lorenz}}
\end{myfig}

Solusi-solusinya cenderung menuju sebuah \emph{atraktor} di ruang, yang disebut
\emph{\myindex{atraktor Lorenz}}.  Dalam hal ini pengamatan dengan strobo tidak
diperlukan; solusi akan cenderung menuju himpunan atraktor.  Sekali lagi kita
tidak dapat benar-benar melihat atraktornya sendiri, tetapi jika kita mencoba
mengikuti suatu solusi cukup lama, seperti pada gambar, kita memperoleh
gambaran yang cukup baik mengenai bentuk atraktor tersebut.  Atraktor Lorenz
juga merupakan atraktor aneh dan memiliki struktur fraktal yang rumit.  Sama
seperti pada persamaan Duffing, yang hendak kita gambar bukan seluruh lintasan.
Kita mulai menggambar lintasan setelah beberapa waktu berlalu, ketika lintasan
sudah dekat dengan atraktor.

Jalur lintasan bukan sekadar bentuk angka delapan yang berulang.  Lintasan
berputar beberapa kali di sebelah kiri dengan jumlah yang tampak acak, kemudian
berputar beberapa kali di sebelah kanan, dan seterusnya.  Karena sistem ini
muncul dalam peramalan cuaca, kita dapat membayangkannya sebagai beberapa hari
bercuaca hangat yang disusul beberapa hari bercuaca dingin.  Tidak mudah
meramalkan kapan cuaca akan berubah, sebagaimana tidak mudah meramalkan jauh
sebelumnya kapan solusi akan berpindah ke sisi lain.  Lihat
\figurevref{nlin:lorenz-graphx} untuk grafik komponen $x$ dari solusi yang
digambar di atas.  Nilai $x$ negatif bersesuaian dengan \myquote{gelung} kiri,
dan nilai $x$ positif bersesuaian dengan \myquote{gelung} kanan.  Di sisi lain,
meskipun kita tidak dapat meramalkan cuaca, kita \emph{dapat} mengatakan
sesuatu tentang iklim---cuaca akan berada di suatu tempat dekat atraktor.

Sebagian besar matematika yang kita pelajari dalam buku ini cukup klasik dan
telah dipahami dengan baik.  Sebaliknya, chaos, termasuk sistem Lorenz, tetap
menjadi pokok penelitian masa kini.  Selain itu, chaos telah diterapkan bukan
hanya dalam sains, tetapi juga dalam seni.

\begin{myfig}
\capstart
\diffyincludegraphics{width=3in}{width=4.5in}{nlin-lorenz-graphx}{%
Grafik sebuah fungsi pada bidang tx untuk t dari 0 hingga 15 dan x kira-kira
antara minus 18 dan 18.  Fungsi bermula sedikit di atas sumbu t dan membentuk
puncak tajam ke atas, di atas sumbu t, lalu puncak tajam ke bawah, kemudian puncak tajam
ke atas, disusul 6 puncak tajam ke bawah yang mula-mula kecil lalu membesar,
kemudian dua puncak tajam ke atas, satu puncak tajam ke bawah, dua puncak tajam
ke atas, satu puncak tajam ke bawah, satu puncak tajam ke atas, dua puncak tajam
ke bawah, lalu dua puncak tajam ke atas.  Puncak-puncak itu berbeda ukuran dan
sedikit berbeda bentuk; tidak tampak pola yang berarti.}
\caption{Grafik komponen $x(t)$ dari solusi.
\label{nlin:lorenz-graphx}}
\end{myfig}

\subsection{Latihan}

\begin{exercise}
Untuk persamaan nonkaotik
$x''+2p x' + \omega_0^2 x = \frac{F_0}{m} \cos (\omega t)$, andaikan kita
melakukan pengamatan dengan strobo berfrekuensi $\omega$ seperti yang telah
disebutkan di atas.  Gunakan solusi periodik tunak yang telah diketahui untuk
menentukan secara tepat titik yang menjadi atraktor bagi penampang Poincar\'e.
\end{exercise}

\begin{samepage}
\begin{exercise}[proyek]
Sebuah atraktor fraktal sederhana dapat digambar melalui permainan chaos
berikut.  Gambarlah tiga titik sudut suatu segitiga dan beri nama, misalnya
$p_1$, $p_2$, dan $p_3$.  Gambarlah suatu titik acak $p$ (titik ini tidak harus
merupakan salah satu dari ketiga titik tadi).  Lempar sebuah dadu untuk memilih
salah satu di antara $p_1$, $p_2$, atau $p_3$ secara acak (misalnya angka 1 dan
4 berarti $p_1$, angka 2 dan 5 berarti $p_2$, serta angka 3 dan 6 berarti
$p_3$).  Andaikan yang terpilih adalah $p_2$.  Misalkan $p_{\text{new}}$ adalah
titik tepat di tengah antara $p$ dan $p_2$.  Gambarlah titik ini, lalu sekarang
gunakan $p$ untuk menyatakan titik baru $p_{\text{new}}$ tersebut.  Ulangi
langkah-langkah ini terus-menerus.  Usahakan bekerja dengan teliti dan gambarlah
sebanyak mungkin iterasi.  Titik-titik Anda akan tertarik menuju apa yang
disebut \emph{\myindex{segitiga Sierpinski}}.  Permainan ini dijalankan dengan
komputer selama 10.000 iterasi untuk menghasilkan gambar dalam
\figurevref{nlin:sierpinski}.
\end{exercise}
\end{samepage}

\begin{myfig}
\capstart
\diffyincludegraphics{width=3in}{width=4.5in}{nlin-sierpinski}{%
Diagram yang disebut segitiga Sierpinski, tersusun dari titik-titik pada
bidang xy.  Titik sudut segitiga berada di (0,0), (1,0), dan (0.5,0.8).
Di dalamnya terdapat segitiga kosong yang lebih kecil, sebangun, dan terbalik,
dengan titik-titik tengah sisi segitiga besar sebagai titik sudutnya.  Tiga
segitiga tegak yang tersisa tidak terisi, tetapi masing-masing memiliki
segitiga kosong seperti yang di tengah; ketiga segitiga tegak tersebut
mengulangi pola ini tanpa akhir.  Dengan demikian, seluruh segitiga tampak
seolah-olah tersusun atas tak terhingga banyak segitiga kosong yang terbalik.}
\caption{10.000 iterasi permainan chaos yang menghasilkan
segitiga Sierpinski. \label{nlin:sierpinski}}
\end{myfig}

\begin{exercise}[proyek]
Buatlah bandul ganda yang dijelaskan dalam teks menggunakan seutas tali dan
dua mur (atau manik-manik berat).  Cobalah berbagai posisi mur tengah dan,
jika diinginkan, gunakan mur dengan bobot berbeda.  Uraikan apa yang Anda
temukan.
\end{exercise}

\begin{exercise}[proyek komputer]
Gunakan perangkat lunak komputer (seperti Matlab, Octave, atau bahkan mungkin
lembar sebar) untuk menggambar solusi persamaan Duffing paksa yang diberikan
dengan metode Euler.  Gambarlah solusi untuk $t$ dari 0 hingga 100 dengan
beberapa ukuran langkah (kecil) yang berbeda.  Bahaslah hasilnya.
\end{exercise}

\setcounter{exercise}{100}

\begin{exercise}
Carilah titik-titik kritis sistem Lorenz beserta linearisasi yang terkait.
\end{exercise}
\exsol{%
Titik-titik kritis: $(0,0,0)$, $(6\sqrt{2},6 \sqrt{2}, 27)$,
$(-6 \sqrt{2},-6 \sqrt{2}, 27)$.
Linearisasi di $(0,0,0)$ dengan $u=x$, $v=y$, $w=z$ adalah
$u' = -10u+10v$, $v'=28u-v$, $w'=-(\nicefrac{8}{3})w$.
Linearisasi di $(6 \sqrt{2},6\sqrt{2},27)$ dengan $u=x-6\sqrt{2}$,
$v=y-6\sqrt{2}$, $w=z-27$ adalah
$u' = -10u+10v$, $v'=u-v-6\sqrt{2}w$, $w'=6\sqrt{2}u+6\sqrt{2}v-(\nicefrac{8}{3})w$.
Linearisasi di $(-6 \sqrt{2},-6\sqrt{2},27)$ dengan $u=x+6\sqrt{2}$,
$v=y+6\sqrt{2}$, $w=z-27$ adalah
$u' = -10u+10v$, $v'=u-v+6\sqrt{2}w$, $w'=-6\sqrt{2}u-6\sqrt{2}v-(\nicefrac{8}{3})w$.
}
